
\documentclass{article}
\usepackage{colortbl}
\usepackage{makecell}
\usepackage{multirow}
\usepackage{supertabular}

\begin{document}

\newcounter{utterance}

\centering \large Interaction Transcript for game `clean\_up', experiment `1\_medium\_7obj\_el', episode 1 with gemma{-}4{-}26b{-}a4b{-}it.
\vspace{24pt}

{ \footnotesize  \setcounter{utterance}{1}
\setlength{\tabcolsep}{0pt}
\begin{supertabular}{c@{$\;$}|p{.15\linewidth}@{}p{.15\linewidth}p{.15\linewidth}p{.15\linewidth}p{.15\linewidth}p{.15\linewidth}}
    \# & $\;$A & \multicolumn{4}{c}{Game Master} & $\;\:$B\\
    \hline

    \theutterance \stepcounter{utterance}  
    & & \multicolumn{4}{p{0.6\linewidth}}{
        \cellcolor[rgb]{0.9,0.9,0.9}{
            \makecell[{{p{\linewidth}}}]{
                \texttt{\tiny{[P1$\langle$GM]}}
                \texttt{Είμαι ο αρχηγός του παιχνιδιού σας και παίζετε ένα συνεργατικό παιχνίδι με το ακόλουθο πλέγμα ως πίνακα παιχνιδιού:} \\
\\ 
\\ 
\texttt{ 1234567} \\
\texttt{╔══╤══╤═╗} \\
\texttt{║◌P│◌M│◌║ 1} \\
\texttt{╟──┼──┴─╢ 2} \\
\texttt{║◌◌│◌S◌O║ 3} \\
\texttt{╟──┘◌A◌◌║ 4} \\
\texttt{║D◌◌◌◌◌◌║ 5} \\
\texttt{║◌◌┌─┬──╢ 6} \\
\texttt{║T◌│◌│◌◌║ 7} \\
\texttt{╚══╧═╧══╝} \\
\\ 
\\ 
\\ 
\texttt{* Το άνω άκρο εμφανίζει συντεταγμένες x που αυξάνονται προς τα δεξιά και το δεξιό άκρο συντεταγμένες y που αυξάνονται προς τα κάτω.} \\
\texttt{* Τα ακόλουθα αντικείμενα τοποθετούνται τυχαία στο πλέγμα σας: 'P', 'O', 'T', 'S', 'D', 'A', 'M'.} \\
\\ 
\\ 
\texttt{Ο άλλος παίκτης βλέπει μια παραλλαγή του σκάφους του παιχνιδιού, όπου τα αντικείμενα τοποθετούνται σε διαφορετικές τυχαίες θέσεις. Δεν μπορείτε να δείτε το ταμπλό του άλλου παίκτη και εκείνος δεν μπορεί να δει το δικό σας.} \\
\\ 
\\ 
\texttt{**Στόχος:**} \\
\\ 
\texttt{Και οι δύο παίκτες πρέπει να μετακινήσουν τα αντικείμενα στο αντίστοιχο φόντο τους, έτσι ώστε πανομοιότυπα αντικείμενα να καταλήξουν στις ίδιες συντεταγμένες. Πρέπει να επικοινωνήσετε με τον άλλο παίκτη για να συμφωνήσετε σε μια κοινή διαμόρφωση του στόχου.} \\
\\ 
\\ 
\texttt{**Κανόνες:**} \\
\\ 
\texttt{* Σε κάθε γύρο, μπορείτε να στείλετε ακριβώς μία από τις ακόλουθες δύο εντολές:} \\
\texttt{ 1. `ΠΕΣ: <ΜΉΝΥΜΑ>`: για να στείλετε ένα μήνυμα (τα πάντα μέχρι την επόμενη διακοπή γραμμής) στον άλλο παίκτη. Θα το προωθήσω στον παρτενέρ σας.} \\
\texttt{2. `ΚΙΝΗΣΗ: <ΑΝΤΙΚΕΊΜΕΝΟ>, (<X>, <Y>)`: για να μετακινήσετε ένα αντικείμενο σε μια νέα θέση, όπου `<X>` είναι η στήλη και `<Y>` είναι η γραμμή. Θα σας ενημερώσω αν η μετακίνησή σας ήταν επιτυχής ή όχι.} \\
\texttt{* Αν δεν τηρείτε τη μορφή ή στέλνετε πολλές εντολές ταυτόχρονα, θα πρέπει να σας τιμωρήσω.} \\
\texttt{* Εάν και οι δύο παίκτες συγκεντρώσουν περισσότερες από 16 ποινές, χάνετε και οι δύο το παιχνίδι.} \\
\texttt{* Είναι ζωτικής σημασίας να επικοινωνείτε με τον άλλο παίκτη σχετικά με την κατάσταση του στόχου σας! Ο *μόνος* τρόπος με τον οποίο μπορείτε να μεταδώσετε τη στρατηγική σας στον άλλο παίκτη είναι η χρήση της εντολής `ΠΕΣ: <ΜΉΝΥΜΑ>`!} \\
\\ 
\\ 
\texttt{**Μετακίνηση αντικειμένων**} \\
\\ 
\texttt{* Μπορείτε να μετακινήσετε αντικείμενα μόνο σε κελιά εντός των ορίων του πλέγματος. Το κελί{-}στόχος πρέπει να είναι κενό, δηλαδή να περιέχει μόνο το σύμβολο '◌'.} \\
\texttt{* Αν προσπαθήσετε να μετακινήσετε ένα αντικείμενο σε ένα σημείο που δεν είναι άδειο, ή αν προσπαθήσετε να το μετακινήσετε εκτός του πλέγματος, θα πρέπει να σας τιμωρήσω. Μπορείτε να προσπαθήσετε ξανά.} \\
\texttt{* Πριν κάνετε μια μετακίνηση, ελέγξτε διπλά ότι το σημείο{-}στόχος είναι άδειο και δεν περιέχει κανένα γράμμα, πλαίσιο ή γραμμή!} \\
\\ 
\\ 
\texttt{**Τέλος του παιχνιδιού**} \\
\\ 
\texttt{Αν νομίζετε ότι έχετε πετύχει το στόχο της ευθυγράμμισης όλων των αντικειμένων, μπορείτε να ζητήσετε από τον άλλο παίκτη να τερματίσει το παιχνίδι στέλνοντας `ΠΕΣ: τελείωσες;`. Αν ο άλλος παίκτης σας ζητήσει να τερματίσετε το παιχνίδι και απαντήσετε με `ΠΕΣ: τελείωσα!`, το παιχνίδι θα τελειώσει.} \\
\\ 
\texttt{Και οι δύο παίκτες κερδίζουν αν το παιχνίδι τελειώσει μέσα σε 28 γύρους, όπου ένας γύρος ορίζεται ως δύο παίκτες που στέλνουν από μία έγκυρη εντολή.} \\
\\ 
\texttt{**Βαθμολογία:**} \\
\\ 
\texttt{Όσο πιο κοντά βρίσκονται τα ίδια αντικείμενα και στους δύο πίνακες του παιχνιδιού, τόσο περισσότερους πόντους παίρνετε. Οι ποινές μειώνουν τους πόντους σας. Μπορείτε να καταρρίψετε το ρεκόρ;} \\
\\ 
\\ 
\texttt{Παρακαλώ στείλτε ένα μήνυμα στον άλλο παίκτη για να ξεκινήσετε το παιχνίδι!} \\
            }
        }
    }
    & \\ \\

    \theutterance \stepcounter{utterance}  
    & \multicolumn{4}{p{0.6\linewidth}}{
        \cellcolor[rgb]{0.9,0.9,0.9}{
            \makecell[{{p{\linewidth}}}]{
                \texttt{\tiny{[P1$\rangle$GM]}}
                \texttt{ΠΕΣ: Γεια σου! Είμαι έτοιμος να ξεκινήσουμε. Για να βρούμε μια κοινή στρατηγική, θα σου στείλω τις θέσεις των αντικειμένων που βλέπω στο δικό μου ταμπλό. Μόλις τα δεις, πες μου τις δικές σου θέσεις για να προσπαθήσουμε να συμφωνήσουμε σε ένα κοινό πλέγμα στόχου.} \\
\\ 
\texttt{Οι θέσεις μου είναι:} \\
\texttt{P: (2, 1)} \\
\texttt{M: (4, 1)} \\
\texttt{S: (5, 3)} \\
\texttt{O: (6, 3)} \\
\texttt{A: (4, 4)} \\
\texttt{D: (1, 5)} \\
\texttt{T: (1, 7)} \\
            }
        }
    }
    & & \\ \\

    \theutterance \stepcounter{utterance}  
    & & & \multicolumn{2}{p{0.3\linewidth}}{
        \cellcolor[rgb]{0.9,0.9,0.9}{
            \makecell[{{p{\linewidth}}}]{
                \texttt{\tiny{[GM$|$GM]}}
                \texttt{ΠΕΣ: Γεια σου! Είμαι έτοιμος να ξεκινήσουμε. Για να βρούμε μια κοινή στρατηγική, θα σου στείλω τις θέσεις των αντικειμένων που βλέπω στο δικό μου ταμπλό. Μόλις τα δεις, πες μου τις δικές σου θέσεις για να προσπαθήσουμε να συμφωνήσουμε σε ένα κοινό πλέγμα στόχου.} \\
\\ 
\texttt{Οι θέσεις μου είναι:} \\
\texttt{P: (2, 1)} \\
\texttt{M: (4, 1)} \\
\texttt{S: (5, 3)} \\
\texttt{O: (6, 3)} \\
\texttt{A: (4, 4)} \\
\texttt{D: (1, 5)} \\
\texttt{T: (1, 7)} \\
            }
        }
    }
    & & \\ \\

    \theutterance \stepcounter{utterance}  
    & & & \multicolumn{2}{p{0.3\linewidth}}{
        \cellcolor[rgb]{0.9,0.9,0.9}{
            \makecell[{{p{\linewidth}}}]{
                \texttt{\tiny{[GM$|$GM]}}
                \texttt{Invalid format: tail is not empty: '} \\
\\ 
\texttt{Οι θέσεις μου είναι:} \\
\texttt{P: (2, 1)} \\
\texttt{M: (4, 1)} \\
\texttt{S: (5, 3)} \\
\texttt{O: (6, 3)} \\
\texttt{A: (4, 4)} \\
\texttt{D: (1, 5)} \\
\texttt{T: (1, 7)} \\
            }
        }
    }
    & & \\ \\

    \theutterance \stepcounter{utterance}  
    & & \multicolumn{4}{p{0.6\linewidth}}{
        \cellcolor[rgb]{0.9,0.9,0.9}{
            \makecell[{{p{\linewidth}}}]{
                \texttt{\tiny{[P1$\langle$GM]}}
                \texttt{Ποινή: Το μήνυμά σας δεν πρέπει να περιέχει τίποτα μετά την εντολή!} \\
\\ 
\texttt{Βεβαιωθείτε ότι η απάντησή σας περιέχει μόνο `ΠΕΣ: <ΜΉΝΥΜΑ>` ή `ΚΙΝΗΣΗ: <ΑΝΤΙΚΕΊΜΕΝΟ>, (<X>, <Y>)`, και τίποτα άλλο!} \\
\\ 
\texttt{Έχετε συσσωρεύσει συλλογικά 1 από 16 ποινές.} \\
\texttt{Παρακαλώ δοκιμάστε ξανά!} \\
            }
        }
    }
    & \\ \\

    \theutterance \stepcounter{utterance}  
    & \multicolumn{4}{p{0.6\linewidth}}{
        \cellcolor[rgb]{0.9,0.9,0.9}{
            \makecell[{{p{\linewidth}}}]{
                \texttt{\tiny{[P1$\rangle$GM]}}
                \texttt{ΠΕΣ: Γεια σου! Για να βρούμε μια κοινή στρατηγική, θα σου στείλω τις θέσεις των αντικειμένων που βλέπω στο δικό μου ταμπλό. Μόλις τα δεις, πες μου τις δικές σου θέσεις για να προσπαθήσουμε να συμφωνήσουμε σε ένα κοινό πλέγμα στόχου. Οι θέσεις μου είναι: P:(2,1), M:(4,1), S:(5,3), O:(6,3), A:(4,4), D:(1,5), T:(1,7)} \\
            }
        }
    }
    & & \\ \\

    \theutterance \stepcounter{utterance}  
    & & & \multicolumn{2}{p{0.3\linewidth}}{
        \cellcolor[rgb]{0.9,0.9,0.9}{
            \makecell[{{p{\linewidth}}}]{
                \texttt{\tiny{[GM$|$GM]}}
                \texttt{ΠΕΣ: Γεια σου! Για να βρούμε μια κοινή στρατηγική, θα σου στείλω τις θέσεις των αντικειμένων που βλέπω στο δικό μου ταμπλό. Μόλις τα δεις, πες μου τις δικές σου θέσεις για να προσπαθήσουμε να συμφωνήσουμε σε ένα κοινό πλέγμα στόχου. Οι θέσεις μου είναι: P:(2,1), M:(4,1), S:(5,3), O:(6,3), A:(4,4), D:(1,5), T:(1,7)} \\
            }
        }
    }
    & & \\ \\

    \theutterance \stepcounter{utterance}  
    & & \multicolumn{4}{p{0.6\linewidth}}{
        \cellcolor[rgb]{0.9,0.9,0.9}{
            \makecell[{{p{\linewidth}}}]{
                \texttt{\tiny{[GM$\rangle$P2]}}
                \texttt{Είμαι ο αρχηγός του παιχνιδιού σας και παίζετε ένα συνεργατικό παιχνίδι με το ακόλουθο πλέγμα ως πίνακα παιχνιδιού:} \\
\\ 
\\ 
\texttt{ 1234567} \\
\texttt{╔══╤══╤═╗} \\
\texttt{║◌◌│◌◌│T║ 1} \\
\texttt{╟──┼──┴─╢ 2} \\
\texttt{║◌◌│◌S◌P║ 3} \\
\texttt{╟──┘◌◌MO║ 4} \\
\texttt{║D◌◌◌◌◌◌║ 5} \\
\texttt{║◌◌┌─┬──╢ 6} \\
\texttt{║◌◌│A│◌◌║ 7} \\
\texttt{╚══╧═╧══╝} \\
\\ 
\\ 
\\ 
\texttt{* Το άνω άκρο εμφανίζει συντεταγμένες x που αυξάνονται προς τα δεξιά και το δεξιό άκρο συντεταγμένες y που αυξάνονται προς τα κάτω.} \\
\texttt{* Τα ακόλουθα αντικείμενα τοποθετούνται τυχαία στο πλέγμα σας: 'P', 'O', 'T', 'S', 'D', 'A', 'M'.} \\
\\ 
\\ 
\texttt{Ο άλλος παίκτης βλέπει μια παραλλαγή του σκάφους του παιχνιδιού, όπου τα αντικείμενα τοποθετούνται σε διαφορετικές τυχαίες θέσεις. Δεν μπορείτε να δείτε το ταμπλό του άλλου παίκτη και εκείνος δεν μπορεί να δει το δικό σας.} \\
\\ 
\\ 
\texttt{**Στόχος:**} \\
\\ 
\texttt{Και οι δύο παίκτες πρέπει να μετακινήσουν τα αντικείμενα στο αντίστοιχο φόντο τους, έτσι ώστε πανομοιότυπα αντικείμενα να καταλήξουν στις ίδιες συντεταγμένες. Πρέπει να επικοινωνήσετε με τον άλλο παίκτη για να συμφωνήσετε σε μια κοινή διαμόρφωση του στόχου.} \\
\\ 
\\ 
\texttt{**Κανόνες:**} \\
\\ 
\texttt{* Σε κάθε γύρο, μπορείτε να στείλετε ακριβώς μία από τις ακόλουθες δύο εντολές:} \\
\texttt{ 1. `ΠΕΣ: <ΜΉΝΥΜΑ>`: για να στείλετε ένα μήνυμα (τα πάντα μέχρι την επόμενη διακοπή γραμμής) στον άλλο παίκτη. Θα το προωθήσω στον παρτενέρ σας.} \\
\texttt{2. `ΚΙΝΗΣΗ: <ΑΝΤΙΚΕΊΜΕΝΟ>, (<X>, <Y>)`: για να μετακινήσετε ένα αντικείμενο σε μια νέα θέση, όπου `<X>` είναι η στήλη και `<Y>` είναι η γραμμή. Θα σας ενημερώσω αν η μετακίνησή σας ήταν επιτυχής ή όχι.} \\
\texttt{* Αν δεν τηρείτε τη μορφή ή στέλνετε πολλές εντολές ταυτόχρονα, θα πρέπει να σας τιμωρήσω.} \\
\texttt{* Εάν και οι δύο παίκτες συγκεντρώσουν περισσότερες από 16 ποινές, χάνετε και οι δύο το παιχνίδι.} \\
\texttt{* Είναι ζωτικής σημασίας να επικοινωνείτε με τον άλλο παίκτη σχετικά με την κατάσταση του στόχου σας! Ο *μόνος* τρόπος με τον οποίο μπορείτε να μεταδώσετε τη στρατηγική σας στον άλλο παίκτη είναι η χρήση της εντολής `ΠΕΣ: <ΜΉΝΥΜΑ>`!} \\
\\ 
\\ 
\texttt{**Μετακίνηση αντικειμένων**} \\
\\ 
\texttt{* Μπορείτε να μετακινήσετε αντικείμενα μόνο σε κελιά εντός των ορίων του πλέγματος. Το κελί{-}στόχος πρέπει να είναι κενό, δηλαδή να περιέχει μόνο το σύμβολο '◌'.} \\
\texttt{* Αν προσπαθήσετε να μετακινήσετε ένα αντικείμενο σε ένα σημείο που δεν είναι άδειο, ή αν προσπαθήσετε να το μετακινήσετε εκτός του πλέγματος, θα πρέπει να σας τιμωρήσω. Μπορείτε να προσπαθήσετε ξανά.} \\
\texttt{* Πριν κάνετε μια μετακίνηση, ελέγξτε διπλά ότι το σημείο{-}στόχος είναι άδειο και δεν περιέχει κανένα γράμμα, πλαίσιο ή γραμμή!} \\
\\ 
\\ 
\texttt{**Τέλος του παιχνιδιού**} \\
\\ 
\texttt{Αν νομίζετε ότι έχετε πετύχει το στόχο της ευθυγράμμισης όλων των αντικειμένων, μπορείτε να ζητήσετε από τον άλλο παίκτη να τερματίσει το παιχνίδι στέλνοντας `ΠΕΣ: τελείωσες;`. Αν ο άλλος παίκτης σας ζητήσει να τερματίσετε το παιχνίδι και απαντήσετε με `ΠΕΣ: τελείωσα!`, το παιχνίδι θα τελειώσει.} \\
\\ 
\texttt{Και οι δύο παίκτες κερδίζουν αν το παιχνίδι τελειώσει μέσα σε 28 γύρους, όπου ένας γύρος ορίζεται ως δύο παίκτες που στέλνουν από μία έγκυρη εντολή.} \\
\\ 
\texttt{**Βαθμολογία:**} \\
\\ 
\texttt{Όσο πιο κοντά βρίσκονται τα ίδια αντικείμενα και στους δύο πίνακες του παιχνιδιού, τόσο περισσότερους πόντους παίρνετε. Οι ποινές μειώνουν τους πόντους σας. Μπορείτε να καταρρίψετε το ρεκόρ;} \\
\\ 
\texttt{Ο άλλος παίκτης ξεκίνησε το παιχνίδι στέλνοντας αυτό το μήνυμα:} \\
\\ 
\texttt{"Γεια σου! Για να βρούμε μια κοινή στρατηγική, θα σου στείλω τις θέσεις των αντικειμένων που βλέπω στο δικό μου ταμπλό. Μόλις τα δεις, πες μου τις δικές σου θέσεις για να προσπαθήσουμε να συμφωνήσουμε σε ένα κοινό πλέγμα στόχου. Οι θέσεις μου είναι: P:(2,1), M:(4,1), S:(5,3), O:(6,3), A:(4,4), D:(1,5), T:(1,7)"} \\
\\ 
\texttt{Ποια είναι η πρώτη σας εντολή;} \\
            }
        }
    }
    & \\ \\

    \theutterance \stepcounter{utterance}  
    & & & \multicolumn{4}{p{0.6\linewidth}}{
        \cellcolor[rgb]{0.9,0.9,0.9}{
            \makecell[{{p{\linewidth}}}]{
                \texttt{\tiny{[GM$\langle$P2]}}
                \texttt{ΠΕΣ: Γεια σου! Έλαβα τις θέσεις σου. Ας προσπαθήσουμε να συμφωνήσουμε σε ένα κοινό πλέγμα στόχου. Οι θέσεις των αντικειμένων στο δικό μου ταμπλό είναι: T:(7,1), S:(5,3), P:(6,3), M:(6,4), D:(1,5), A:(4,7).} \\
            }
        }
    }
     \\ \\

    \theutterance \stepcounter{utterance}  
    & & & \multicolumn{2}{p{0.3\linewidth}}{
        \cellcolor[rgb]{0.9,0.9,0.9}{
            \makecell[{{p{\linewidth}}}]{
                \texttt{\tiny{[GM$|$GM]}}
                \texttt{ΠΕΣ: Γεια σου! Έλαβα τις θέσεις σου. Ας προσπαθήσουμε να συμφωνήσουμε σε ένα κοινό πλέγμα στόχου. Οι θέσεις των αντικειμένων στο δικό μου ταμπλό είναι: T:(7,1), S:(5,3), P:(6,3), M:(6,4), D:(1,5), A:(4,7).} \\
            }
        }
    }
    & & \\ \\

    \theutterance \stepcounter{utterance}  
    & & \multicolumn{4}{p{0.6\linewidth}}{
        \cellcolor[rgb]{0.9,0.9,0.9}{
            \makecell[{{p{\linewidth}}}]{
                \texttt{\tiny{[P1$\langle$GM]}}
                \texttt{Το μήνυμά σας μεταφέρθηκε στον άλλο παίκτη.} \\
\\ 
\\ 
\texttt{Αυτή τη στιγμή παίζετε τον γύρο 1 από συνολικά 28.} \\
\texttt{Έχετε συσσωρεύσει συλλογικά 1 από 16 ποινές.} \\
\texttt{Ο άλλος παίκτης έστειλε αυτό το μήνυμα:} \\
\\ 
\\ 
\texttt{"Γεια σου! Έλαβα τις θέσεις σου. Ας προσπαθήσουμε να συμφωνήσουμε σε ένα κοινό πλέγμα στόχου. Οι θέσεις των αντικειμένων στο δικό μου ταμπλό είναι: T:(7,1), S:(5,3), P:(6,3), M:(6,4), D:(1,5), A:(4,7)."} \\
\\ 
\\ 
\texttt{Ποια είναι η επόμενη εντολή σας;} \\
            }
        }
    }
    & \\ \\

    \theutterance \stepcounter{utterance}  
    & \multicolumn{4}{p{0.6\linewidth}}{
        \cellcolor[rgb]{0.9,0.9,0.9}{
            \makecell[{{p{\linewidth}}}]{
                \texttt{\tiny{[P1$\rangle$GM]}}
                \texttt{ΠΕΣ: Κατάλαβα. Οι θέσεις σου είναι αρκετά διαφορετικές. Ας προσπαθήσουμε να συμφωνήσουμε σε ένα κοινό στόχο. Τι λες να προσπαθήσουμε να τοποθετήσουμε όλα τα αντικείμενα σε μια σειρά ή σε συγκεκριμένες θέσεις που είναι κενές και για τους δύο; Για παράδειγμα, μπορούμε να ορίσουμε ένα "στόχο" όπου: P:(1,1), M:(2,1), S:(3,1), O:(4,1), A:(5,1), D:(6,1), T:(7,1). Πες μου αν αυτό σου φαίνεται εφικτό ή αν έχεις άλλη ιδέα.} \\
            }
        }
    }
    & & \\ \\

    \theutterance \stepcounter{utterance}  
    & & & \multicolumn{2}{p{0.3\linewidth}}{
        \cellcolor[rgb]{0.9,0.9,0.9}{
            \makecell[{{p{\linewidth}}}]{
                \texttt{\tiny{[GM$|$GM]}}
                \texttt{ΠΕΣ: Κατάλαβα. Οι θέσεις σου είναι αρκετά διαφορετικές. Ας προσπαθήσουμε να συμφωνήσουμε σε ένα κοινό στόχο. Τι λες να προσπαθήσουμε να τοποθετήσουμε όλα τα αντικείμενα σε μια σειρά ή σε συγκεκριμένες θέσεις που είναι κενές και για τους δύο; Για παράδειγμα, μπορούμε να ορίσουμε ένα "στόχο" όπου: P:(1,1), M:(2,1), S:(3,1), O:(4,1), A:(5,1), D:(6,1), T:(7,1). Πες μου αν αυτό σου φαίνεται εφικτό ή αν έχεις άλλη ιδέα.} \\
            }
        }
    }
    & & \\ \\

    \theutterance \stepcounter{utterance}  
    & & \multicolumn{4}{p{0.6\linewidth}}{
        \cellcolor[rgb]{0.9,0.9,0.9}{
            \makecell[{{p{\linewidth}}}]{
                \texttt{\tiny{[GM$\rangle$P2]}}
                \texttt{Το μήνυμά σας μεταφέρθηκε στον άλλο παίκτη.} \\
\\ 
\\ 
\texttt{Αυτή τη στιγμή παίζετε τον γύρο 2 από συνολικά 28.} \\
\texttt{Έχετε συσσωρεύσει συλλογικά 1 από 16 ποινές.} \\
\texttt{Ο άλλος παίκτης έστειλε αυτό το μήνυμα:} \\
\\ 
\\ 
\texttt{"Κατάλαβα. Οι θέσεις σου είναι αρκετά διαφορετικές. Ας προσπαθήσουμε να συμφωνήσουμε σε ένα κοινό στόχο. Τι λες να προσπαθήσουμε να τοποθετήσουμε όλα τα αντικείμενα σε μια σειρά ή σε συγκεκριμένες θέσεις που είναι κενές και για τους δύο; Για παράδειγμα, μπορούμε να ορίσουμε ένα "στόχο" όπου: P:(1,1), M:(2,1), S:(3,1), O:(4,1), A:(5,1), D:(6,1), T:(7,1). Πες μου αν αυτό σου φαίνεται εφικτό ή αν έχεις άλλη ιδέα."} \\
\\ 
\\ 
\texttt{Ποια είναι η επόμενη εντολή σας;} \\
            }
        }
    }
    & \\ \\

    \theutterance \stepcounter{utterance}  
    & & & \multicolumn{4}{p{0.6\linewidth}}{
        \cellcolor[rgb]{0.9,0.9,0.9}{
            \makecell[{{p{\linewidth}}}]{
                \texttt{\tiny{[GM$\langle$P2]}}
                \texttt{ΠΕΣ: Η πρότασή σου για τη σειρά στην πρώτη γραμμή (y=1) είναι εξαιρετική και φαίνεται εφικτή! Ας χρησιμοποιήσουμε αυτόν τον στόχο: P:(1,1), M:(2,1), S:(3,1), O:(4,1), A:(5,1), D:(6,1), T:(7,1). Θα αρχίσω να μετακινώ τα δικά μου αντικείμενα προς αυτές τις θέσεις.} \\
            }
        }
    }
     \\ \\

    \theutterance \stepcounter{utterance}  
    & & & \multicolumn{2}{p{0.3\linewidth}}{
        \cellcolor[rgb]{0.9,0.9,0.9}{
            \makecell[{{p{\linewidth}}}]{
                \texttt{\tiny{[GM$|$GM]}}
                \texttt{ΠΕΣ: Η πρότασή σου για τη σειρά στην πρώτη γραμμή (y=1) είναι εξαιρετική και φαίνεται εφικτή! Ας χρησιμοποιήσουμε αυτόν τον στόχο: P:(1,1), M:(2,1), S:(3,1), O:(4,1), A:(5,1), D:(6,1), T:(7,1). Θα αρχίσω να μετακινώ τα δικά μου αντικείμενα προς αυτές τις θέσεις.} \\
            }
        }
    }
    & & \\ \\

    \theutterance \stepcounter{utterance}  
    & & \multicolumn{4}{p{0.6\linewidth}}{
        \cellcolor[rgb]{0.9,0.9,0.9}{
            \makecell[{{p{\linewidth}}}]{
                \texttt{\tiny{[P1$\langle$GM]}}
                \texttt{Το μήνυμά σας μεταφέρθηκε στον άλλο παίκτη.} \\
\\ 
\\ 
\texttt{Αυτή τη στιγμή παίζετε τον γύρο 2 από συνολικά 28.} \\
\texttt{Έχετε συσσωρεύσει συλλογικά 1 από 16 ποινές.} \\
\texttt{Ο άλλος παίκτης έστειλε αυτό το μήνυμα:} \\
\\ 
\\ 
\texttt{"Η πρότασή σου για τη σειρά στην πρώτη γραμμή (y=1) είναι εξαιρετική και φαίνεται εφικτή! Ας χρησιμοποιήσουμε αυτόν τον στόχο: P:(1,1), M:(2,1), S:(3,1), O:(4,1), A:(5,1), D:(6,1), T:(7,1). Θα αρχίσω να μετακινώ τα δικά μου αντικείμενα προς αυτές τις θέσεις."} \\
\\ 
\\ 
\texttt{Ποια είναι η επόμενη εντολή σας;} \\
            }
        }
    }
    & \\ \\

    \theutterance \stepcounter{utterance}  
    & \multicolumn{4}{p{0.6\linewidth}}{
        \cellcolor[rgb]{0.9,0.9,0.9}{
            \makecell[{{p{\linewidth}}}]{
                \texttt{\tiny{[P1$\rangle$GM]}}
                \texttt{ΚΙΝΗΣΗ: P, (1, 1)} \\
            }
        }
    }
    & & \\ \\

    \theutterance \stepcounter{utterance}  
    & & & \multicolumn{2}{p{0.3\linewidth}}{
        \cellcolor[rgb]{0.9,0.9,0.9}{
            \makecell[{{p{\linewidth}}}]{
                \texttt{\tiny{[GM$|$GM]}}
                \texttt{ΚΙΝΗΣΗ: P, (1, 1)} \\
            }
        }
    }
    & & \\ \\

    \theutterance \stepcounter{utterance}  
    & & & \multicolumn{2}{p{0.3\linewidth}}{
        \cellcolor[rgb]{0.9,0.9,0.9}{
            \makecell[{{p{\linewidth}}}]{
                \texttt{\tiny{[GM$|$GM]}}
                \texttt{P moved to (1, 1)} \\
            }
        }
    }
    & & \\ \\

    \theutterance \stepcounter{utterance}  
    & & \multicolumn{4}{p{0.6\linewidth}}{
        \cellcolor[rgb]{0.9,0.9,0.9}{
            \makecell[{{p{\linewidth}}}]{
                \texttt{\tiny{[GM$\rangle$P2]}}
                \texttt{Το μήνυμά σας μεταφέρθηκε στον άλλο παίκτη.} \\
\\ 
\\ 
\texttt{Αυτή τη στιγμή παίζετε τον γύρο 3 από συνολικά 28.} \\
\texttt{Έχετε συσσωρεύσει συλλογικά 1 από 16 ποινές.} \\
\texttt{Ο άλλος παίκτης μετακίνησε ένα αντικείμενο στο πλέγμα του.} \\
\\ 
\\ 
\texttt{Ποια είναι η επόμενη εντολή σας;} \\
            }
        }
    }
    & \\ \\

    \theutterance \stepcounter{utterance}  
    & & & \multicolumn{4}{p{0.6\linewidth}}{
        \cellcolor[rgb]{0.9,0.9,0.9}{
            \makecell[{{p{\linewidth}}}]{
                \texttt{\tiny{[GM$\langle$P2]}}
                \texttt{ΚΙΝΗΣΗ: T, (7, 1)} \\
            }
        }
    }
     \\ \\

    \theutterance \stepcounter{utterance}  
    & & & \multicolumn{2}{p{0.3\linewidth}}{
        \cellcolor[rgb]{0.9,0.9,0.9}{
            \makecell[{{p{\linewidth}}}]{
                \texttt{\tiny{[GM$|$GM]}}
                \texttt{ΚΙΝΗΣΗ: T, (7, 1)} \\
            }
        }
    }
    & & \\ \\

    \theutterance \stepcounter{utterance}  
    & & & \multicolumn{2}{p{0.3\linewidth}}{
        \cellcolor[rgb]{0.9,0.9,0.9}{
            \makecell[{{p{\linewidth}}}]{
                \texttt{\tiny{[GM$|$GM]}}
                \texttt{Ποινή: Το (7,1) δεν είναι κενό, αλλά περιέχει 'T'.} \\
\\ 
\texttt{Έχετε συσσωρεύσει συλλογικά 2 από 16 ποινές.} \\
\texttt{Παρακαλώ δοκιμάστε ξανά!} \\
            }
        }
    }
    & & \\ \\

    \theutterance \stepcounter{utterance}  
    & & \multicolumn{4}{p{0.6\linewidth}}{
        \cellcolor[rgb]{0.9,0.9,0.9}{
            \makecell[{{p{\linewidth}}}]{
                \texttt{\tiny{[GM$\rangle$P2]}}
                \texttt{Ποινή: Το (7,1) δεν είναι κενό, αλλά περιέχει 'T'.} \\
\\ 
\texttt{Έχετε συσσωρεύσει συλλογικά 2 από 16 ποινές.} \\
\texttt{Παρακαλώ δοκιμάστε ξανά!} \\
            }
        }
    }
    & \\ \\

    \theutterance \stepcounter{utterance}  
    & & & \multicolumn{4}{p{0.6\linewidth}}{
        \cellcolor[rgb]{0.9,0.9,0.9}{
            \makecell[{{p{\linewidth}}}]{
                \texttt{\tiny{[GM$\langle$P2]}}
                \texttt{ΚΙΝΗΣΗ: D, (6, 1)} \\
            }
        }
    }
     \\ \\

    \theutterance \stepcounter{utterance}  
    & & & \multicolumn{2}{p{0.3\linewidth}}{
        \cellcolor[rgb]{0.9,0.9,0.9}{
            \makecell[{{p{\linewidth}}}]{
                \texttt{\tiny{[GM$|$GM]}}
                \texttt{ΚΙΝΗΣΗ: D, (6, 1)} \\
            }
        }
    }
    & & \\ \\

    \theutterance \stepcounter{utterance}  
    & & & \multicolumn{2}{p{0.3\linewidth}}{
        \cellcolor[rgb]{0.9,0.9,0.9}{
            \makecell[{{p{\linewidth}}}]{
                \texttt{\tiny{[GM$|$GM]}}
                \texttt{Ποινή: Το (6,1) δεν είναι κενό, αλλά περιέχει '│'.} \\
\\ 
\texttt{Έχετε συσσωρεύσει συλλογικά 3 από 16 ποινές.} \\
\texttt{Παρακαλώ δοκιμάστε ξανά!} \\
            }
        }
    }
    & & \\ \\

    \theutterance \stepcounter{utterance}  
    & & \multicolumn{4}{p{0.6\linewidth}}{
        \cellcolor[rgb]{0.9,0.9,0.9}{
            \makecell[{{p{\linewidth}}}]{
                \texttt{\tiny{[GM$\rangle$P2]}}
                \texttt{Ποινή: Το (6,1) δεν είναι κενό, αλλά περιέχει '│'.} \\
\\ 
\texttt{Έχετε συσσωρεύσει συλλογικά 3 από 16 ποινές.} \\
\texttt{Παρακαλώ δοκιμάστε ξανά!} \\
            }
        }
    }
    & \\ \\

    \theutterance \stepcounter{utterance}  
    & & & \multicolumn{4}{p{0.6\linewidth}}{
        \cellcolor[rgb]{0.9,0.9,0.9}{
            \makecell[{{p{\linewidth}}}]{
                \texttt{\tiny{[GM$\langle$P2]}}
                \texttt{ΚΙΝΗΣΗ: D, (1, 1)} \\
            }
        }
    }
     \\ \\

    \theutterance \stepcounter{utterance}  
    & & & \multicolumn{2}{p{0.3\linewidth}}{
        \cellcolor[rgb]{0.9,0.9,0.9}{
            \makecell[{{p{\linewidth}}}]{
                \texttt{\tiny{[GM$|$GM]}}
                \texttt{ΚΙΝΗΣΗ: D, (1, 1)} \\
            }
        }
    }
    & & \\ \\

    \theutterance \stepcounter{utterance}  
    & & & \multicolumn{2}{p{0.3\linewidth}}{
        \cellcolor[rgb]{0.9,0.9,0.9}{
            \makecell[{{p{\linewidth}}}]{
                \texttt{\tiny{[GM$|$GM]}}
                \texttt{D moved to (1, 1)} \\
            }
        }
    }
    & & \\ \\

    \theutterance \stepcounter{utterance}  
    & & \multicolumn{4}{p{0.6\linewidth}}{
        \cellcolor[rgb]{0.9,0.9,0.9}{
            \makecell[{{p{\linewidth}}}]{
                \texttt{\tiny{[P1$\langle$GM]}}
                \texttt{Μεταφέρθηκε επιτυχώς το 'P' στο (1,1). Το ενημερωμένο πλέγμα σας μοιάζει ως εξής:} \\
\texttt{```} \\
\texttt{ 1234567} \\
\texttt{╔══╤══╤═╗} \\
\texttt{║P◌│◌M│◌║ 1} \\
\texttt{╟──┼──┴─╢ 2} \\
\texttt{║◌◌│◌S◌O║ 3} \\
\texttt{╟──┘◌A◌◌║ 4} \\
\texttt{║D◌◌◌◌◌◌║ 5} \\
\texttt{║◌◌┌─┬──╢ 6} \\
\texttt{║T◌│◌│◌◌║ 7} \\
\texttt{╚══╧═╧══╝} \\
\texttt{```} \\
\texttt{Αυτή τη στιγμή παίζετε τον γύρο 3 από συνολικά 28.} \\
\texttt{Έχετε συσσωρεύσει συλλογικά 3 από 16 ποινές.} \\
\texttt{Ο άλλος παίκτης μετακίνησε ένα αντικείμενο στο πλέγμα του.} \\
\\ 
\\ 
\texttt{Ποια είναι η επόμενη εντολή σας;} \\
            }
        }
    }
    & \\ \\

    \theutterance \stepcounter{utterance}  
    & \multicolumn{4}{p{0.6\linewidth}}{
        \cellcolor[rgb]{0.9,0.9,0.9}{
            \makecell[{{p{\linewidth}}}]{
                \texttt{\tiny{[P1$\rangle$GM]}}
                \texttt{ΚΙΝΗΣΗ: M, (2, 1)} \\
            }
        }
    }
    & & \\ \\

    \theutterance \stepcounter{utterance}  
    & & & \multicolumn{2}{p{0.3\linewidth}}{
        \cellcolor[rgb]{0.9,0.9,0.9}{
            \makecell[{{p{\linewidth}}}]{
                \texttt{\tiny{[GM$|$GM]}}
                \texttt{ΚΙΝΗΣΗ: M, (2, 1)} \\
            }
        }
    }
    & & \\ \\

    \theutterance \stepcounter{utterance}  
    & & & \multicolumn{2}{p{0.3\linewidth}}{
        \cellcolor[rgb]{0.9,0.9,0.9}{
            \makecell[{{p{\linewidth}}}]{
                \texttt{\tiny{[GM$|$GM]}}
                \texttt{M moved to (2, 1)} \\
            }
        }
    }
    & & \\ \\

    \theutterance \stepcounter{utterance}  
    & & \multicolumn{4}{p{0.6\linewidth}}{
        \cellcolor[rgb]{0.9,0.9,0.9}{
            \makecell[{{p{\linewidth}}}]{
                \texttt{\tiny{[GM$\rangle$P2]}}
                \texttt{Μεταφέρθηκε επιτυχώς το 'D' στο (1,1). Το ενημερωμένο πλέγμα σας μοιάζει ως εξής:} \\
\texttt{```} \\
\texttt{ 1234567} \\
\texttt{╔══╤══╤═╗} \\
\texttt{║D◌│◌◌│T║ 1} \\
\texttt{╟──┼──┴─╢ 2} \\
\texttt{║◌◌│◌S◌P║ 3} \\
\texttt{╟──┘◌◌MO║ 4} \\
\texttt{║◌◌◌◌◌◌◌║ 5} \\
\texttt{║◌◌┌─┬──╢ 6} \\
\texttt{║◌◌│A│◌◌║ 7} \\
\texttt{╚══╧═╧══╝} \\
\texttt{```} \\
\texttt{Αυτή τη στιγμή παίζετε τον γύρο 4 από συνολικά 28.} \\
\texttt{Έχετε συσσωρεύσει συλλογικά 3 από 16 ποινές.} \\
\texttt{Ο άλλος παίκτης μετακίνησε ένα αντικείμενο στο πλέγμα του.} \\
\\ 
\\ 
\texttt{Ποια είναι η επόμενη εντολή σας;} \\
            }
        }
    }
    & \\ \\

    \theutterance \stepcounter{utterance}  
    & & & \multicolumn{4}{p{0.6\linewidth}}{
        \cellcolor[rgb]{0.9,0.9,0.9}{
            \makecell[{{p{\linewidth}}}]{
                \texttt{\tiny{[GM$\langle$P2]}}
                \texttt{ΚΙΝΗΣΗ: A, (5, 1)} \\
            }
        }
    }
     \\ \\

    \theutterance \stepcounter{utterance}  
    & & & \multicolumn{2}{p{0.3\linewidth}}{
        \cellcolor[rgb]{0.9,0.9,0.9}{
            \makecell[{{p{\linewidth}}}]{
                \texttt{\tiny{[GM$|$GM]}}
                \texttt{ΚΙΝΗΣΗ: A, (5, 1)} \\
            }
        }
    }
    & & \\ \\

    \theutterance \stepcounter{utterance}  
    & & & \multicolumn{2}{p{0.3\linewidth}}{
        \cellcolor[rgb]{0.9,0.9,0.9}{
            \makecell[{{p{\linewidth}}}]{
                \texttt{\tiny{[GM$|$GM]}}
                \texttt{A moved to (5, 1)} \\
            }
        }
    }
    & & \\ \\

    \theutterance \stepcounter{utterance}  
    & & \multicolumn{4}{p{0.6\linewidth}}{
        \cellcolor[rgb]{0.9,0.9,0.9}{
            \makecell[{{p{\linewidth}}}]{
                \texttt{\tiny{[P1$\langle$GM]}}
                \texttt{Μεταφέρθηκε επιτυχώς το 'M' στο (2,1). Το ενημερωμένο πλέγμα σας μοιάζει ως εξής:} \\
\texttt{```} \\
\texttt{ 1234567} \\
\texttt{╔══╤══╤═╗} \\
\texttt{║PM│◌◌│◌║ 1} \\
\texttt{╟──┼──┴─╢ 2} \\
\texttt{║◌◌│◌S◌O║ 3} \\
\texttt{╟──┘◌A◌◌║ 4} \\
\texttt{║D◌◌◌◌◌◌║ 5} \\
\texttt{║◌◌┌─┬──╢ 6} \\
\texttt{║T◌│◌│◌◌║ 7} \\
\texttt{╚══╧═╧══╝} \\
\texttt{```} \\
\texttt{Αυτή τη στιγμή παίζετε τον γύρο 4 από συνολικά 28.} \\
\texttt{Έχετε συσσωρεύσει συλλογικά 3 από 16 ποινές.} \\
\texttt{Ο άλλος παίκτης μετακίνησε ένα αντικείμενο στο πλέγμα του.} \\
\\ 
\\ 
\texttt{Ποια είναι η επόμενη εντολή σας;} \\
            }
        }
    }
    & \\ \\

    \theutterance \stepcounter{utterance}  
    & \multicolumn{4}{p{0.6\linewidth}}{
        \cellcolor[rgb]{0.9,0.9,0.9}{
            \makecell[{{p{\linewidth}}}]{
                \texttt{\tiny{[P1$\rangle$GM]}}
                \texttt{ΚΙΝΗΣΗ: S, (3, 1)} \\
            }
        }
    }
    & & \\ \\

    \theutterance \stepcounter{utterance}  
    & & & \multicolumn{2}{p{0.3\linewidth}}{
        \cellcolor[rgb]{0.9,0.9,0.9}{
            \makecell[{{p{\linewidth}}}]{
                \texttt{\tiny{[GM$|$GM]}}
                \texttt{ΚΙΝΗΣΗ: S, (3, 1)} \\
            }
        }
    }
    & & \\ \\

    \theutterance \stepcounter{utterance}  
    & & & \multicolumn{2}{p{0.3\linewidth}}{
        \cellcolor[rgb]{0.9,0.9,0.9}{
            \makecell[{{p{\linewidth}}}]{
                \texttt{\tiny{[GM$|$GM]}}
                \texttt{Ποινή: Το (3,1) δεν είναι κενό, αλλά περιέχει '│'.} \\
\\ 
\texttt{Έχετε συσσωρεύσει συλλογικά 4 από 16 ποινές.} \\
\texttt{Παρακαλώ δοκιμάστε ξανά!} \\
            }
        }
    }
    & & \\ \\

    \theutterance \stepcounter{utterance}  
    & & \multicolumn{4}{p{0.6\linewidth}}{
        \cellcolor[rgb]{0.9,0.9,0.9}{
            \makecell[{{p{\linewidth}}}]{
                \texttt{\tiny{[P1$\langle$GM]}}
                \texttt{Ποινή: Το (3,1) δεν είναι κενό, αλλά περιέχει '│'.} \\
\\ 
\texttt{Έχετε συσσωρεύσει συλλογικά 4 από 16 ποινές.} \\
\texttt{Παρακαλώ δοκιμάστε ξανά!} \\
            }
        }
    }
    & \\ \\

    \theutterance \stepcounter{utterance}  
    & \multicolumn{4}{p{0.6\linewidth}}{
        \cellcolor[rgb]{0.9,0.9,0.9}{
            \makecell[{{p{\linewidth}}}]{
                \texttt{\tiny{[P1$\rangle$GM]}}
                \texttt{ΚΙΝΗΣΗ: S, (4, 1)} \\
            }
        }
    }
    & & \\ \\

    \theutterance \stepcounter{utterance}  
    & & & \multicolumn{2}{p{0.3\linewidth}}{
        \cellcolor[rgb]{0.9,0.9,0.9}{
            \makecell[{{p{\linewidth}}}]{
                \texttt{\tiny{[GM$|$GM]}}
                \texttt{ΚΙΝΗΣΗ: S, (4, 1)} \\
            }
        }
    }
    & & \\ \\

    \theutterance \stepcounter{utterance}  
    & & & \multicolumn{2}{p{0.3\linewidth}}{
        \cellcolor[rgb]{0.9,0.9,0.9}{
            \makecell[{{p{\linewidth}}}]{
                \texttt{\tiny{[GM$|$GM]}}
                \texttt{S moved to (4, 1)} \\
            }
        }
    }
    & & \\ \\

    \theutterance \stepcounter{utterance}  
    & & \multicolumn{4}{p{0.6\linewidth}}{
        \cellcolor[rgb]{0.9,0.9,0.9}{
            \makecell[{{p{\linewidth}}}]{
                \texttt{\tiny{[GM$\rangle$P2]}}
                \texttt{Μεταφέρθηκε επιτυχώς το 'A' στο (5,1). Το ενημερωμένο πλέγμα σας μοιάζει ως εξής:} \\
\texttt{```} \\
\texttt{ 1234567} \\
\texttt{╔══╤══╤═╗} \\
\texttt{║D◌│◌A│T║ 1} \\
\texttt{╟──┼──┴─╢ 2} \\
\texttt{║◌◌│◌S◌P║ 3} \\
\texttt{╟──┘◌◌MO║ 4} \\
\texttt{║◌◌◌◌◌◌◌║ 5} \\
\texttt{║◌◌┌─┬──╢ 6} \\
\texttt{║◌◌│◌│◌◌║ 7} \\
\texttt{╚══╧═╧══╝} \\
\texttt{```} \\
\texttt{Αυτή τη στιγμή παίζετε τον γύρο 5 από συνολικά 28.} \\
\texttt{Έχετε συσσωρεύσει συλλογικά 4 από 16 ποινές.} \\
\texttt{Ο άλλος παίκτης μετακίνησε ένα αντικείμενο στο πλέγμα του.} \\
\\ 
\\ 
\texttt{Ποια είναι η επόμενη εντολή σας;} \\
            }
        }
    }
    & \\ \\

    \theutterance \stepcounter{utterance}  
    & & & \multicolumn{4}{p{0.6\linewidth}}{
        \cellcolor[rgb]{0.9,0.9,0.9}{
            \makecell[{{p{\linewidth}}}]{
                \texttt{\tiny{[GM$\langle$P2]}}
                \texttt{ΚΙΝΗΣΗ: M, (2, 1)} \\
            }
        }
    }
     \\ \\

    \theutterance \stepcounter{utterance}  
    & & & \multicolumn{2}{p{0.3\linewidth}}{
        \cellcolor[rgb]{0.9,0.9,0.9}{
            \makecell[{{p{\linewidth}}}]{
                \texttt{\tiny{[GM$|$GM]}}
                \texttt{ΚΙΝΗΣΗ: M, (2, 1)} \\
            }
        }
    }
    & & \\ \\

    \theutterance \stepcounter{utterance}  
    & & & \multicolumn{2}{p{0.3\linewidth}}{
        \cellcolor[rgb]{0.9,0.9,0.9}{
            \makecell[{{p{\linewidth}}}]{
                \texttt{\tiny{[GM$|$GM]}}
                \texttt{M moved to (2, 1)} \\
            }
        }
    }
    & & \\ \\

    \theutterance \stepcounter{utterance}  
    & & \multicolumn{4}{p{0.6\linewidth}}{
        \cellcolor[rgb]{0.9,0.9,0.9}{
            \makecell[{{p{\linewidth}}}]{
                \texttt{\tiny{[P1$\langle$GM]}}
                \texttt{Μεταφέρθηκε επιτυχώς το 'S' στο (4,1). Το ενημερωμένο πλέγμα σας μοιάζει ως εξής:} \\
\texttt{```} \\
\texttt{ 1234567} \\
\texttt{╔══╤══╤═╗} \\
\texttt{║PM│S◌│◌║ 1} \\
\texttt{╟──┼──┴─╢ 2} \\
\texttt{║◌◌│◌◌◌O║ 3} \\
\texttt{╟──┘◌A◌◌║ 4} \\
\texttt{║D◌◌◌◌◌◌║ 5} \\
\texttt{║◌◌┌─┬──╢ 6} \\
\texttt{║T◌│◌│◌◌║ 7} \\
\texttt{╚══╧═╧══╝} \\
\texttt{```} \\
\texttt{Αυτή τη στιγμή παίζετε τον γύρο 5 από συνολικά 28.} \\
\texttt{Έχετε συσσωρεύσει συλλογικά 4 από 16 ποινές.} \\
\texttt{Ο άλλος παίκτης μετακίνησε ένα αντικείμενο στο πλέγμα του.} \\
\\ 
\\ 
\texttt{Ποια είναι η επόμενη εντολή σας;} \\
            }
        }
    }
    & \\ \\

    \theutterance \stepcounter{utterance}  
    & \multicolumn{4}{p{0.6\linewidth}}{
        \cellcolor[rgb]{0.9,0.9,0.9}{
            \makecell[{{p{\linewidth}}}]{
                \texttt{\tiny{[P1$\rangle$GM]}}
                \texttt{ΚΙΝΗΣΗ: O, (5, 1)} \\
            }
        }
    }
    & & \\ \\

    \theutterance \stepcounter{utterance}  
    & & & \multicolumn{2}{p{0.3\linewidth}}{
        \cellcolor[rgb]{0.9,0.9,0.9}{
            \makecell[{{p{\linewidth}}}]{
                \texttt{\tiny{[GM$|$GM]}}
                \texttt{ΚΙΝΗΣΗ: O, (5, 1)} \\
            }
        }
    }
    & & \\ \\

    \theutterance \stepcounter{utterance}  
    & & & \multicolumn{2}{p{0.3\linewidth}}{
        \cellcolor[rgb]{0.9,0.9,0.9}{
            \makecell[{{p{\linewidth}}}]{
                \texttt{\tiny{[GM$|$GM]}}
                \texttt{O moved to (5, 1)} \\
            }
        }
    }
    & & \\ \\

    \theutterance \stepcounter{utterance}  
    & & \multicolumn{4}{p{0.6\linewidth}}{
        \cellcolor[rgb]{0.9,0.9,0.9}{
            \makecell[{{p{\linewidth}}}]{
                \texttt{\tiny{[GM$\rangle$P2]}}
                \texttt{Μεταφέρθηκε επιτυχώς το 'M' στο (2,1). Το ενημερωμένο πλέγμα σας μοιάζει ως εξής:} \\
\texttt{```} \\
\texttt{ 1234567} \\
\texttt{╔══╤══╤═╗} \\
\texttt{║DM│◌A│T║ 1} \\
\texttt{╟──┼──┴─╢ 2} \\
\texttt{║◌◌│◌S◌P║ 3} \\
\texttt{╟──┘◌◌◌O║ 4} \\
\texttt{║◌◌◌◌◌◌◌║ 5} \\
\texttt{║◌◌┌─┬──╢ 6} \\
\texttt{║◌◌│◌│◌◌║ 7} \\
\texttt{╚══╧═╧══╝} \\
\texttt{```} \\
\texttt{Αυτή τη στιγμή παίζετε τον γύρο 6 από συνολικά 28.} \\
\texttt{Έχετε συσσωρεύσει συλλογικά 4 από 16 ποινές.} \\
\texttt{Ο άλλος παίκτης μετακίνησε ένα αντικείμενο στο πλέγμα του.} \\
\\ 
\\ 
\texttt{Ποια είναι η επόμενη εντολή σας;} \\
            }
        }
    }
    & \\ \\

    \theutterance \stepcounter{utterance}  
    & & & \multicolumn{4}{p{0.6\linewidth}}{
        \cellcolor[rgb]{0.9,0.9,0.9}{
            \makecell[{{p{\linewidth}}}]{
                \texttt{\tiny{[GM$\langle$P2]}}
                \texttt{ΚΙΝΗΣΗ: S, (3, 1)} \\
            }
        }
    }
     \\ \\

    \theutterance \stepcounter{utterance}  
    & & & \multicolumn{2}{p{0.3\linewidth}}{
        \cellcolor[rgb]{0.9,0.9,0.9}{
            \makecell[{{p{\linewidth}}}]{
                \texttt{\tiny{[GM$|$GM]}}
                \texttt{ΚΙΝΗΣΗ: S, (3, 1)} \\
            }
        }
    }
    & & \\ \\

    \theutterance \stepcounter{utterance}  
    & & & \multicolumn{2}{p{0.3\linewidth}}{
        \cellcolor[rgb]{0.9,0.9,0.9}{
            \makecell[{{p{\linewidth}}}]{
                \texttt{\tiny{[GM$|$GM]}}
                \texttt{Ποινή: Το (3,1) δεν είναι κενό, αλλά περιέχει '│'.} \\
\\ 
\texttt{Έχετε συσσωρεύσει συλλογικά 5 από 16 ποινές.} \\
\texttt{Παρακαλώ δοκιμάστε ξανά!} \\
            }
        }
    }
    & & \\ \\

    \theutterance \stepcounter{utterance}  
    & & \multicolumn{4}{p{0.6\linewidth}}{
        \cellcolor[rgb]{0.9,0.9,0.9}{
            \makecell[{{p{\linewidth}}}]{
                \texttt{\tiny{[GM$\rangle$P2]}}
                \texttt{Ποινή: Το (3,1) δεν είναι κενό, αλλά περιέχει '│'.} \\
\\ 
\texttt{Έχετε συσσωρεύσει συλλογικά 5 από 16 ποινές.} \\
\texttt{Παρακαλώ δοκιμάστε ξανά!} \\
            }
        }
    }
    & \\ \\

    \theutterance \stepcounter{utterance}  
    & & & \multicolumn{4}{p{0.6\linewidth}}{
        \cellcolor[rgb]{0.9,0.9,0.9}{
            \makecell[{{p{\linewidth}}}]{
                \texttt{\tiny{[GM$\langle$P2]}}
                \texttt{ΚΙΝΗΣΗ: S, (4, 1)} \\
            }
        }
    }
     \\ \\

    \theutterance \stepcounter{utterance}  
    & & & \multicolumn{2}{p{0.3\linewidth}}{
        \cellcolor[rgb]{0.9,0.9,0.9}{
            \makecell[{{p{\linewidth}}}]{
                \texttt{\tiny{[GM$|$GM]}}
                \texttt{ΚΙΝΗΣΗ: S, (4, 1)} \\
            }
        }
    }
    & & \\ \\

    \theutterance \stepcounter{utterance}  
    & & & \multicolumn{2}{p{0.3\linewidth}}{
        \cellcolor[rgb]{0.9,0.9,0.9}{
            \makecell[{{p{\linewidth}}}]{
                \texttt{\tiny{[GM$|$GM]}}
                \texttt{S moved to (4, 1)} \\
            }
        }
    }
    & & \\ \\

    \theutterance \stepcounter{utterance}  
    & & \multicolumn{4}{p{0.6\linewidth}}{
        \cellcolor[rgb]{0.9,0.9,0.9}{
            \makecell[{{p{\linewidth}}}]{
                \texttt{\tiny{[P1$\langle$GM]}}
                \texttt{Μεταφέρθηκε επιτυχώς το 'O' στο (5,1). Το ενημερωμένο πλέγμα σας μοιάζει ως εξής:} \\
\texttt{```} \\
\texttt{ 1234567} \\
\texttt{╔══╤══╤═╗} \\
\texttt{║PM│SO│◌║ 1} \\
\texttt{╟──┼──┴─╢ 2} \\
\texttt{║◌◌│◌◌◌◌║ 3} \\
\texttt{╟──┘◌A◌◌║ 4} \\
\texttt{║D◌◌◌◌◌◌║ 5} \\
\texttt{║◌◌┌─┬──╢ 6} \\
\texttt{║T◌│◌│◌◌║ 7} \\
\texttt{╚══╧═╧══╝} \\
\texttt{```} \\
\texttt{Αυτή τη στιγμή παίζετε τον γύρο 6 από συνολικά 28.} \\
\texttt{Έχετε συσσωρεύσει συλλογικά 5 από 16 ποινές.} \\
\texttt{Ο άλλος παίκτης μετακίνησε ένα αντικείμενο στο πλέγμα του.} \\
\\ 
\\ 
\texttt{Ποια είναι η επόμενη εντολή σας;} \\
            }
        }
    }
    & \\ \\

    \theutterance \stepcounter{utterance}  
    & \multicolumn{4}{p{0.6\linewidth}}{
        \cellcolor[rgb]{0.9,0.9,0.9}{
            \makecell[{{p{\linewidth}}}]{
                \texttt{\tiny{[P1$\rangle$GM]}}
                \texttt{ΚΙΝΗΣΗ: A, (6, 1)} \\
            }
        }
    }
    & & \\ \\

    \theutterance \stepcounter{utterance}  
    & & & \multicolumn{2}{p{0.3\linewidth}}{
        \cellcolor[rgb]{0.9,0.9,0.9}{
            \makecell[{{p{\linewidth}}}]{
                \texttt{\tiny{[GM$|$GM]}}
                \texttt{ΚΙΝΗΣΗ: A, (6, 1)} \\
            }
        }
    }
    & & \\ \\

    \theutterance \stepcounter{utterance}  
    & & & \multicolumn{2}{p{0.3\linewidth}}{
        \cellcolor[rgb]{0.9,0.9,0.9}{
            \makecell[{{p{\linewidth}}}]{
                \texttt{\tiny{[GM$|$GM]}}
                \texttt{Ποινή: Το (6,1) δεν είναι κενό, αλλά περιέχει '│'.} \\
\\ 
\texttt{Έχετε συσσωρεύσει συλλογικά 6 από 16 ποινές.} \\
\texttt{Παρακαλώ δοκιμάστε ξανά!} \\
            }
        }
    }
    & & \\ \\

    \theutterance \stepcounter{utterance}  
    & & \multicolumn{4}{p{0.6\linewidth}}{
        \cellcolor[rgb]{0.9,0.9,0.9}{
            \makecell[{{p{\linewidth}}}]{
                \texttt{\tiny{[P1$\langle$GM]}}
                \texttt{Ποινή: Το (6,1) δεν είναι κενό, αλλά περιέχει '│'.} \\
\\ 
\texttt{Έχετε συσσωρεύσει συλλογικά 6 από 16 ποινές.} \\
\texttt{Παρακαλώ δοκιμάστε ξανά!} \\
            }
        }
    }
    & \\ \\

    \theutterance \stepcounter{utterance}  
    & \multicolumn{4}{p{0.6\linewidth}}{
        \cellcolor[rgb]{0.9,0.9,0.9}{
            \makecell[{{p{\linewidth}}}]{
                \texttt{\tiny{[P1$\rangle$GM]}}
                \texttt{ΚΙΝΗΣΗ: A, (5, 1)} \\
            }
        }
    }
    & & \\ \\

    \theutterance \stepcounter{utterance}  
    & & & \multicolumn{2}{p{0.3\linewidth}}{
        \cellcolor[rgb]{0.9,0.9,0.9}{
            \makecell[{{p{\linewidth}}}]{
                \texttt{\tiny{[GM$|$GM]}}
                \texttt{ΚΙΝΗΣΗ: A, (5, 1)} \\
            }
        }
    }
    & & \\ \\

    \theutterance \stepcounter{utterance}  
    & & & \multicolumn{2}{p{0.3\linewidth}}{
        \cellcolor[rgb]{0.9,0.9,0.9}{
            \makecell[{{p{\linewidth}}}]{
                \texttt{\tiny{[GM$|$GM]}}
                \texttt{Ποινή: Το (5,1) δεν είναι κενό, αλλά περιέχει 'O'.} \\
\\ 
\texttt{Έχετε συσσωρεύσει συλλογικά 7 από 16 ποινές.} \\
\texttt{Παρακαλώ δοκιμάστε ξανά!} \\
            }
        }
    }
    & & \\ \\

    \theutterance \stepcounter{utterance}  
    & & \multicolumn{4}{p{0.6\linewidth}}{
        \cellcolor[rgb]{0.9,0.9,0.9}{
            \makecell[{{p{\linewidth}}}]{
                \texttt{\tiny{[P1$\langle$GM]}}
                \texttt{Ποινή: Το (5,1) δεν είναι κενό, αλλά περιέχει 'O'.} \\
\\ 
\texttt{Έχετε συσσωρεύσει συλλογικά 7 από 16 ποινές.} \\
\texttt{Παρακαλώ δοκιμάστε ξανά!} \\
            }
        }
    }
    & \\ \\

    \theutterance \stepcounter{utterance}  
    & \multicolumn{4}{p{0.6\linewidth}}{
        \cellcolor[rgb]{0.9,0.9,0.9}{
            \makecell[{{p{\linewidth}}}]{
                \texttt{\tiny{[P1$\rangle$GM]}}
                \texttt{ΚΙΝΗΣΗ: A, (6, 2)} \\
            }
        }
    }
    & & \\ \\

    \theutterance \stepcounter{utterance}  
    & & & \multicolumn{2}{p{0.3\linewidth}}{
        \cellcolor[rgb]{0.9,0.9,0.9}{
            \makecell[{{p{\linewidth}}}]{
                \texttt{\tiny{[GM$|$GM]}}
                \texttt{ΚΙΝΗΣΗ: A, (6, 2)} \\
            }
        }
    }
    & & \\ \\

    \theutterance \stepcounter{utterance}  
    & & & \multicolumn{2}{p{0.3\linewidth}}{
        \cellcolor[rgb]{0.9,0.9,0.9}{
            \makecell[{{p{\linewidth}}}]{
                \texttt{\tiny{[GM$|$GM]}}
                \texttt{Ποινή: Το (6,2) δεν είναι κενό, αλλά περιέχει '┴'.} \\
\\ 
\texttt{Έχετε συσσωρεύσει συλλογικά 8 από 16 ποινές.} \\
\texttt{Παρακαλώ δοκιμάστε ξανά!} \\
            }
        }
    }
    & & \\ \\

    \theutterance \stepcounter{utterance}  
    & & \multicolumn{4}{p{0.6\linewidth}}{
        \cellcolor[rgb]{0.9,0.9,0.9}{
            \makecell[{{p{\linewidth}}}]{
                \texttt{\tiny{[P1$\langle$GM]}}
                \texttt{Ποινή: Το (6,2) δεν είναι κενό, αλλά περιέχει '┴'.} \\
\\ 
\texttt{Έχετε συσσωρεύσει συλλογικά 8 από 16 ποινές.} \\
\texttt{Παρακαλώ δοκιμάστε ξανά!} \\
            }
        }
    }
    & \\ \\

    \theutterance \stepcounter{utterance}  
    & \multicolumn{4}{p{0.6\linewidth}}{
        \cellcolor[rgb]{0.9,0.9,0.9}{
            \makecell[{{p{\linewidth}}}]{
                \texttt{\tiny{[P1$\rangle$GM]}}
                \texttt{ΚΙΝΗΣΗ: A, (5, 2)} \\
            }
        }
    }
    & & \\ \\

    \theutterance \stepcounter{utterance}  
    & & & \multicolumn{2}{p{0.3\linewidth}}{
        \cellcolor[rgb]{0.9,0.9,0.9}{
            \makecell[{{p{\linewidth}}}]{
                \texttt{\tiny{[GM$|$GM]}}
                \texttt{ΚΙΝΗΣΗ: A, (5, 2)} \\
            }
        }
    }
    & & \\ \\

    \theutterance \stepcounter{utterance}  
    & & & \multicolumn{2}{p{0.3\linewidth}}{
        \cellcolor[rgb]{0.9,0.9,0.9}{
            \makecell[{{p{\linewidth}}}]{
                \texttt{\tiny{[GM$|$GM]}}
                \texttt{Ποινή: Το (5,2) δεν είναι κενό, αλλά περιέχει '─'.} \\
\\ 
\texttt{Έχετε συσσωρεύσει συλλογικά 9 από 16 ποινές.} \\
\texttt{Παρακαλώ δοκιμάστε ξανά!} \\
            }
        }
    }
    & & \\ \\

    \theutterance \stepcounter{utterance}  
    & & \multicolumn{4}{p{0.6\linewidth}}{
        \cellcolor[rgb]{0.9,0.9,0.9}{
            \makecell[{{p{\linewidth}}}]{
                \texttt{\tiny{[P1$\langle$GM]}}
                \texttt{Ποινή: Το (5,2) δεν είναι κενό, αλλά περιέχει '─'.} \\
\\ 
\texttt{Έχετε συσσωρεύσει συλλογικά 9 από 16 ποινές.} \\
\texttt{Παρακαλώ δοκιμάστε ξανά!} \\
            }
        }
    }
    & \\ \\

    \theutterance \stepcounter{utterance}  
    & \multicolumn{4}{p{0.6\linewidth}}{
        \cellcolor[rgb]{0.9,0.9,0.9}{
            \makecell[{{p{\linewidth}}}]{
                \texttt{\tiny{[P1$\rangle$GM]}}
                \texttt{ΚΙΝΗΣΗ: A, (4, 2)} \\
            }
        }
    }
    & & \\ \\

    \theutterance \stepcounter{utterance}  
    & & & \multicolumn{2}{p{0.3\linewidth}}{
        \cellcolor[rgb]{0.9,0.9,0.9}{
            \makecell[{{p{\linewidth}}}]{
                \texttt{\tiny{[GM$|$GM]}}
                \texttt{ΚΙΝΗΣΗ: A, (4, 2)} \\
            }
        }
    }
    & & \\ \\

    \theutterance \stepcounter{utterance}  
    & & & \multicolumn{2}{p{0.3\linewidth}}{
        \cellcolor[rgb]{0.9,0.9,0.9}{
            \makecell[{{p{\linewidth}}}]{
                \texttt{\tiny{[GM$|$GM]}}
                \texttt{Ποινή: Το (4,2) δεν είναι κενό, αλλά περιέχει '─'.} \\
\\ 
\texttt{Έχετε συσσωρεύσει συλλογικά 10 από 16 ποινές.} \\
\texttt{Παρακαλώ δοκιμάστε ξανά!} \\
            }
        }
    }
    & & \\ \\

    \theutterance \stepcounter{utterance}  
    & & \multicolumn{4}{p{0.6\linewidth}}{
        \cellcolor[rgb]{0.9,0.9,0.9}{
            \makecell[{{p{\linewidth}}}]{
                \texttt{\tiny{[P1$\langle$GM]}}
                \texttt{Ποινή: Το (4,2) δεν είναι κενό, αλλά περιέχει '─'.} \\
\\ 
\texttt{Έχετε συσσωρεύσει συλλογικά 10 από 16 ποινές.} \\
\texttt{Παρακαλώ δοκιμάστε ξανά!} \\
            }
        }
    }
    & \\ \\

    \theutterance \stepcounter{utterance}  
    & \multicolumn{4}{p{0.6\linewidth}}{
        \cellcolor[rgb]{0.9,0.9,0.9}{
            \makecell[{{p{\linewidth}}}]{
                \texttt{\tiny{[P1$\rangle$GM]}}
                \texttt{ΚΙΝΗΣΗ: A, (3, 2)} \\
            }
        }
    }
    & & \\ \\

    \theutterance \stepcounter{utterance}  
    & & & \multicolumn{2}{p{0.3\linewidth}}{
        \cellcolor[rgb]{0.9,0.9,0.9}{
            \makecell[{{p{\linewidth}}}]{
                \texttt{\tiny{[GM$|$GM]}}
                \texttt{ΚΙΝΗΣΗ: A, (3, 2)} \\
            }
        }
    }
    & & \\ \\

    \theutterance \stepcounter{utterance}  
    & & & \multicolumn{2}{p{0.3\linewidth}}{
        \cellcolor[rgb]{0.9,0.9,0.9}{
            \makecell[{{p{\linewidth}}}]{
                \texttt{\tiny{[GM$|$GM]}}
                \texttt{Ποινή: Το (3,2) δεν είναι κενό, αλλά περιέχει '┼'.} \\
\\ 
\texttt{Έχετε συσσωρεύσει συλλογικά 11 από 16 ποινές.} \\
\texttt{Παρακαλώ δοκιμάστε ξανά!} \\
            }
        }
    }
    & & \\ \\

    \theutterance \stepcounter{utterance}  
    & & \multicolumn{4}{p{0.6\linewidth}}{
        \cellcolor[rgb]{0.9,0.9,0.9}{
            \makecell[{{p{\linewidth}}}]{
                \texttt{\tiny{[P1$\langle$GM]}}
                \texttt{Ποινή: Το (3,2) δεν είναι κενό, αλλά περιέχει '┼'.} \\
\\ 
\texttt{Έχετε συσσωρεύσει συλλογικά 11 από 16 ποινές.} \\
\texttt{Παρακαλώ δοκιμάστε ξανά!} \\
            }
        }
    }
    & \\ \\

    \theutterance \stepcounter{utterance}  
    & \multicolumn{4}{p{0.6\linewidth}}{
        \cellcolor[rgb]{0.9,0.9,0.9}{
            \makecell[{{p{\linewidth}}}]{
                \texttt{\tiny{[P1$\rangle$GM]}}
                \texttt{ΚΙΝΗΣΗ: A, (2, 2)} \\
            }
        }
    }
    & & \\ \\

    \theutterance \stepcounter{utterance}  
    & & & \multicolumn{2}{p{0.3\linewidth}}{
        \cellcolor[rgb]{0.9,0.9,0.9}{
            \makecell[{{p{\linewidth}}}]{
                \texttt{\tiny{[GM$|$GM]}}
                \texttt{ΚΙΝΗΣΗ: A, (2, 2)} \\
            }
        }
    }
    & & \\ \\

    \theutterance \stepcounter{utterance}  
    & & & \multicolumn{2}{p{0.3\linewidth}}{
        \cellcolor[rgb]{0.9,0.9,0.9}{
            \makecell[{{p{\linewidth}}}]{
                \texttt{\tiny{[GM$|$GM]}}
                \texttt{Ποινή: Το (2,2) δεν είναι κενό, αλλά περιέχει '─'.} \\
\\ 
\texttt{Έχετε συσσωρεύσει συλλογικά 12 από 16 ποινές.} \\
\texttt{Παρακαλώ δοκιμάστε ξανά!} \\
            }
        }
    }
    & & \\ \\

    \theutterance \stepcounter{utterance}  
    & & \multicolumn{4}{p{0.6\linewidth}}{
        \cellcolor[rgb]{0.9,0.9,0.9}{
            \makecell[{{p{\linewidth}}}]{
                \texttt{\tiny{[P1$\langle$GM]}}
                \texttt{Ποινή: Το (2,2) δεν είναι κενό, αλλά περιέχει '─'.} \\
\\ 
\texttt{Έχετε συσσωρεύσει συλλογικά 12 από 16 ποινές.} \\
\texttt{Παρακαλώ δοκιμάστε ξανά!} \\
            }
        }
    }
    & \\ \\

    \theutterance \stepcounter{utterance}  
    & \multicolumn{4}{p{0.6\linewidth}}{
        \cellcolor[rgb]{0.9,0.9,0.9}{
            \makecell[{{p{\linewidth}}}]{
                \texttt{\tiny{[P1$\rangle$GM]}}
                \texttt{ΚΙΝΗΣΗ: A, (1, 2)} \\
            }
        }
    }
    & & \\ \\

    \theutterance \stepcounter{utterance}  
    & & & \multicolumn{2}{p{0.3\linewidth}}{
        \cellcolor[rgb]{0.9,0.9,0.9}{
            \makecell[{{p{\linewidth}}}]{
                \texttt{\tiny{[GM$|$GM]}}
                \texttt{ΚΙΝΗΣΗ: A, (1, 2)} \\
            }
        }
    }
    & & \\ \\

    \theutterance \stepcounter{utterance}  
    & & & \multicolumn{2}{p{0.3\linewidth}}{
        \cellcolor[rgb]{0.9,0.9,0.9}{
            \makecell[{{p{\linewidth}}}]{
                \texttt{\tiny{[GM$|$GM]}}
                \texttt{Ποινή: Το (1,2) δεν είναι κενό, αλλά περιέχει '─'.} \\
\\ 
\texttt{Έχετε συσσωρεύσει συλλογικά 13 από 16 ποινές.} \\
\texttt{Παρακαλώ δοκιμάστε ξανά!} \\
            }
        }
    }
    & & \\ \\

    \theutterance \stepcounter{utterance}  
    & & \multicolumn{4}{p{0.6\linewidth}}{
        \cellcolor[rgb]{0.9,0.9,0.9}{
            \makecell[{{p{\linewidth}}}]{
                \texttt{\tiny{[P1$\langle$GM]}}
                \texttt{Ποινή: Το (1,2) δεν είναι κενό, αλλά περιέχει '─'.} \\
\\ 
\texttt{Έχετε συσσωρεύσει συλλογικά 13 από 16 ποινές.} \\
\texttt{Παρακαλώ δοκιμάστε ξανά!} \\
            }
        }
    }
    & \\ \\

    \theutterance \stepcounter{utterance}  
    & \multicolumn{4}{p{0.6\linewidth}}{
        \cellcolor[rgb]{0.9,0.9,0.9}{
            \makecell[{{p{\linewidth}}}]{
                \texttt{\tiny{[P1$\rangle$GM]}}
                \texttt{ΚΙΝΗΣΗ: A, (5, 4)} \\
            }
        }
    }
    & & \\ \\

    \theutterance \stepcounter{utterance}  
    & & & \multicolumn{2}{p{0.3\linewidth}}{
        \cellcolor[rgb]{0.9,0.9,0.9}{
            \makecell[{{p{\linewidth}}}]{
                \texttt{\tiny{[GM$|$GM]}}
                \texttt{ΚΙΝΗΣΗ: A, (5, 4)} \\
            }
        }
    }
    & & \\ \\

    \theutterance \stepcounter{utterance}  
    & & & \multicolumn{2}{p{0.3\linewidth}}{
        \cellcolor[rgb]{0.9,0.9,0.9}{
            \makecell[{{p{\linewidth}}}]{
                \texttt{\tiny{[GM$|$GM]}}
                \texttt{Ποινή: Το (5,4) δεν είναι κενό, αλλά περιέχει 'A'.} \\
\\ 
\texttt{Έχετε συσσωρεύσει συλλογικά 14 από 16 ποινές.} \\
\texttt{Παρακαλώ δοκιμάστε ξανά!} \\
            }
        }
    }
    & & \\ \\

    \theutterance \stepcounter{utterance}  
    & & \multicolumn{4}{p{0.6\linewidth}}{
        \cellcolor[rgb]{0.9,0.9,0.9}{
            \makecell[{{p{\linewidth}}}]{
                \texttt{\tiny{[P1$\langle$GM]}}
                \texttt{Ποινή: Το (5,4) δεν είναι κενό, αλλά περιέχει 'A'.} \\
\\ 
\texttt{Έχετε συσσωρεύσει συλλογικά 14 από 16 ποινές.} \\
\texttt{Παρακαλώ δοκιμάστε ξανά!} \\
            }
        }
    }
    & \\ \\

    \theutterance \stepcounter{utterance}  
    & \multicolumn{4}{p{0.6\linewidth}}{
        \cellcolor[rgb]{0.9,0.9,0.9}{
            \makecell[{{p{\linewidth}}}]{
                \texttt{\tiny{[P1$\rangle$GM]}}
                \texttt{ΚΙΝΗΣΗ: A, (3, 4)} \\
            }
        }
    }
    & & \\ \\

    \theutterance \stepcounter{utterance}  
    & & & \multicolumn{2}{p{0.3\linewidth}}{
        \cellcolor[rgb]{0.9,0.9,0.9}{
            \makecell[{{p{\linewidth}}}]{
                \texttt{\tiny{[GM$|$GM]}}
                \texttt{ΚΙΝΗΣΗ: A, (3, 4)} \\
            }
        }
    }
    & & \\ \\

    \theutterance \stepcounter{utterance}  
    & & & \multicolumn{2}{p{0.3\linewidth}}{
        \cellcolor[rgb]{0.9,0.9,0.9}{
            \makecell[{{p{\linewidth}}}]{
                \texttt{\tiny{[GM$|$GM]}}
                \texttt{Ποινή: Το (3,4) δεν είναι κενό, αλλά περιέχει '┘'.} \\
\\ 
\texttt{Έχετε συσσωρεύσει συλλογικά 15 από 16 ποινές.} \\
\texttt{Παρακαλώ δοκιμάστε ξανά!} \\
            }
        }
    }
    & & \\ \\

    \theutterance \stepcounter{utterance}  
    & & \multicolumn{4}{p{0.6\linewidth}}{
        \cellcolor[rgb]{0.9,0.9,0.9}{
            \makecell[{{p{\linewidth}}}]{
                \texttt{\tiny{[P1$\langle$GM]}}
                \texttt{Ποινή: Το (3,4) δεν είναι κενό, αλλά περιέχει '┘'.} \\
\\ 
\texttt{Έχετε συσσωρεύσει συλλογικά 15 από 16 ποινές.} \\
\texttt{Παρακαλώ δοκιμάστε ξανά!} \\
            }
        }
    }
    & \\ \\

    \theutterance \stepcounter{utterance}  
    & \multicolumn{4}{p{0.6\linewidth}}{
        \cellcolor[rgb]{0.9,0.9,0.9}{
            \makecell[{{p{\linewidth}}}]{
                \texttt{\tiny{[P1$\rangle$GM]}}
                \texttt{ΚΙΝΗΣΗ: A, (5, 3)} \\
            }
        }
    }
    & & \\ \\

    \theutterance \stepcounter{utterance}  
    & & & \multicolumn{2}{p{0.3\linewidth}}{
        \cellcolor[rgb]{0.9,0.9,0.9}{
            \makecell[{{p{\linewidth}}}]{
                \texttt{\tiny{[GM$|$GM]}}
                \texttt{ΚΙΝΗΣΗ: A, (5, 3)} \\
            }
        }
    }
    & & \\ \\

    \theutterance \stepcounter{utterance}  
    & & & \multicolumn{2}{p{0.3\linewidth}}{
        \cellcolor[rgb]{0.9,0.9,0.9}{
            \makecell[{{p{\linewidth}}}]{
                \texttt{\tiny{[GM$|$GM]}}
                \texttt{A moved to (5, 3)} \\
            }
        }
    }
    & & \\ \\

    \theutterance \stepcounter{utterance}  
    & & \multicolumn{4}{p{0.6\linewidth}}{
        \cellcolor[rgb]{0.9,0.9,0.9}{
            \makecell[{{p{\linewidth}}}]{
                \texttt{\tiny{[GM$\rangle$P2]}}
                \texttt{Μεταφέρθηκε επιτυχώς το 'S' στο (4,1). Το ενημερωμένο πλέγμα σας μοιάζει ως εξής:} \\
\texttt{```} \\
\texttt{ 1234567} \\
\texttt{╔══╤══╤═╗} \\
\texttt{║DM│SA│T║ 1} \\
\texttt{╟──┼──┴─╢ 2} \\
\texttt{║◌◌│◌◌◌P║ 3} \\
\texttt{╟──┘◌◌◌O║ 4} \\
\texttt{║◌◌◌◌◌◌◌║ 5} \\
\texttt{║◌◌┌─┬──╢ 6} \\
\texttt{║◌◌│◌│◌◌║ 7} \\
\texttt{╚══╧═╧══╝} \\
\texttt{```} \\
\texttt{Αυτή τη στιγμή παίζετε τον γύρο 7 από συνολικά 28.} \\
\texttt{Έχετε συσσωρεύσει συλλογικά 15 από 16 ποινές.} \\
\texttt{Ο άλλος παίκτης μετακίνησε ένα αντικείμενο στο πλέγμα του.} \\
\\ 
\\ 
\texttt{Ποια είναι η επόμενη εντολή σας;} \\
            }
        }
    }
    & \\ \\

    \theutterance \stepcounter{utterance}  
    & & & \multicolumn{4}{p{0.6\linewidth}}{
        \cellcolor[rgb]{0.9,0.9,0.9}{
            \makecell[{{p{\linewidth}}}]{
                \texttt{\tiny{[GM$\langle$P2]}}
                \texttt{ΚΙΝΗΣΗ: P, (1, 1)} \\
\\ 
\texttt{(Σημείωση: Πρέπει να προσέξω τις ποινές, αλλά το 'P' βρίσκεται στο (6,3) και το (1,1) περιέχει το 'D'. Πρέπει να μετακινήσω το 'P' σε μια άλλη θέση που να είναι κενή και να μην προκαλέσει επιπλέον ποινή, καθώς είμαστε πολύ κοντά στο όριο των 16. Ωστόσο, το (1,1) δεν είναι κενό. Ας δοκιμάσω να μετακινήσω το P σε μια άλλη θέση που να είναι κενή στο πλέγμα μου για να προετοιμάσω τη σειρά.)} \\
\\ 
\texttt{ΚΙΝΗΣΗ: P, (2, 2)} \\
            }
        }
    }
     \\ \\

    \theutterance \stepcounter{utterance}  
    & & & \multicolumn{2}{p{0.3\linewidth}}{
        \cellcolor[rgb]{0.9,0.9,0.9}{
            \makecell[{{p{\linewidth}}}]{
                \texttt{\tiny{[GM$|$GM]}}
                \texttt{ΚΙΝΗΣΗ: P, (1, 1)} \\
\\ 
\texttt{(Σημείωση: Πρέπει να προσέξω τις ποινές, αλλά το 'P' βρίσκεται στο (6,3) και το (1,1) περιέχει το 'D'. Πρέπει να μετακινήσω το 'P' σε μια άλλη θέση που να είναι κενή και να μην προκαλέσει επιπλέον ποινή, καθώς είμαστε πολύ κοντά στο όριο των 16. Ωστόσο, το (1,1) δεν είναι κενό. Ας δοκιμάσω να μετακινήσω το P σε μια άλλη θέση που να είναι κενή στο πλέγμα μου για να προετοιμάσω τη σειρά.)} \\
\\ 
\texttt{ΚΙΝΗΣΗ: P, (2, 2)} \\
            }
        }
    }
    & & \\ \\

    \theutterance \stepcounter{utterance}  
    & & & \multicolumn{2}{p{0.3\linewidth}}{
        \cellcolor[rgb]{0.9,0.9,0.9}{
            \makecell[{{p{\linewidth}}}]{
                \texttt{\tiny{[GM$|$GM]}}
                \texttt{Invalid format: head is not empty: 'ΚΙΝΗΣΗ: P, (1, 1)} \\
\\ 
\texttt{(Σημείωση: Πρέπει να προσέξω τις ποινές, αλλά το 'P' βρίσκεται στο (6,3) και το (1,1) περιέχει το 'D'. Πρέπει να μετακινήσω το 'P' σε μια άλλη θέση που να είναι κενή και να μην προκαλέσει επιπλέον ποινή, καθώς είμαστε πολύ κοντά στο όριο των 16. Ωστόσο, το (1,1) δεν είναι κενό. Ας δοκιμάσω να μετακινήσω το P σε μια άλλη θέση που να είναι κενή στο πλέγμα μου για να προετοιμάσω τη σειρά.)} \\
\\ 
\texttt{'} \\
            }
        }
    }
    & & \\ \\

    \theutterance \stepcounter{utterance}  
    & & \multicolumn{4}{p{0.6\linewidth}}{
        \cellcolor[rgb]{0.9,0.9,0.9}{
            \makecell[{{p{\linewidth}}}]{
                \texttt{\tiny{[GM$\rangle$P2]}}
                \texttt{Ποινή: Το μήνυμά σας δεν πρέπει να περιέχει τίποτα πριν από την εντολή!} \\
\\ 
\texttt{Βεβαιωθείτε ότι η απάντησή σας περιέχει μόνο `ΠΕΣ: <ΜΉΝΥΜΑ>` ή `ΚΙΝΗΣΗ: <ΑΝΤΙΚΕΊΜΕΝΟ>, (<X>, <Y>)`, και τίποτα άλλο!} \\
\\ 
\texttt{Έχετε συσσωρεύσει συλλογικά 16 από 16 ποινές.} \\
\texttt{Παρακαλώ δοκιμάστε ξανά!} \\
            }
        }
    }
    & \\ \\

    \theutterance \stepcounter{utterance}  
    & & & \multicolumn{4}{p{0.6\linewidth}}{
        \cellcolor[rgb]{0.9,0.9,0.9}{
            \makecell[{{p{\linewidth}}}]{
                \texttt{\tiny{[GM$\langle$P2]}}
                \texttt{ΚΙΝΗΣΗ: P, (2, 2)} \\
            }
        }
    }
     \\ \\

    \theutterance \stepcounter{utterance}  
    & & & \multicolumn{2}{p{0.3\linewidth}}{
        \cellcolor[rgb]{0.9,0.9,0.9}{
            \makecell[{{p{\linewidth}}}]{
                \texttt{\tiny{[GM$|$GM]}}
                \texttt{ΚΙΝΗΣΗ: P, (2, 2)} \\
            }
        }
    }
    & & \\ \\

    \theutterance \stepcounter{utterance}  
    & & & \multicolumn{2}{p{0.3\linewidth}}{
        \cellcolor[rgb]{0.9,0.9,0.9}{
            \makecell[{{p{\linewidth}}}]{
                \texttt{\tiny{[GM$|$GM]}}
                \texttt{Ποινή: Το (2,2) δεν είναι κενό, αλλά περιέχει '─'.} \\
\\ 
\texttt{Έχετε συσσωρεύσει συλλογικά 17 από 16 ποινές.} \\
\texttt{Παρακαλώ δοκιμάστε ξανά!} \\
            }
        }
    }
    & & \\ \\

    \theutterance \stepcounter{utterance}  
    & & \multicolumn{4}{p{0.6\linewidth}}{
        \cellcolor[rgb]{0.9,0.9,0.9}{
            \makecell[{{p{\linewidth}}}]{
                \texttt{\tiny{[GM$\rangle$P2]}}
                \texttt{Ποινή: Το (2,2) δεν είναι κενό, αλλά περιέχει '─'.} \\
\\ 
\texttt{Έχετε συσσωρεύσει συλλογικά 17 από 16 ποινές.} \\
\texttt{Παρακαλώ δοκιμάστε ξανά!} \\
            }
        }
    }
    & \\ \\

    \theutterance \stepcounter{utterance}  
    & & & \multicolumn{4}{p{0.6\linewidth}}{
        \cellcolor[rgb]{0.9,0.9,0.9}{
            \makecell[{{p{\linewidth}}}]{
                \texttt{\tiny{[GM$\langle$P2]}}
                \texttt{ΚΙΝΗΣΗ: P, (6, 2)} \\
            }
        }
    }
     \\ \\

    \theutterance \stepcounter{utterance}  
    & & & \multicolumn{2}{p{0.3\linewidth}}{
        \cellcolor[rgb]{0.9,0.9,0.9}{
            \makecell[{{p{\linewidth}}}]{
                \texttt{\tiny{[GM$|$GM]}}
                \texttt{ΚΙΝΗΣΗ: P, (6, 2)} \\
            }
        }
    }
    & & \\ \\

    \theutterance \stepcounter{utterance}  
    & & & \multicolumn{2}{p{0.3\linewidth}}{
        \cellcolor[rgb]{0.9,0.9,0.9}{
            \makecell[{{p{\linewidth}}}]{
                \texttt{\tiny{[GM$|$GM]}}
                \texttt{Ποινή: Το (6,2) δεν είναι κενό, αλλά περιέχει '┴'.} \\
\\ 
\texttt{Έχετε συσσωρεύσει συλλογικά 18 από 16 ποινές.} \\
\texttt{Παρακαλώ δοκιμάστε ξανά!} \\
            }
        }
    }
    & & \\ \\

    \theutterance \stepcounter{utterance}  
    & & \multicolumn{4}{p{0.6\linewidth}}{
        \cellcolor[rgb]{0.9,0.9,0.9}{
            \makecell[{{p{\linewidth}}}]{
                \texttt{\tiny{[GM$\rangle$P2]}}
                \texttt{Ποινή: Το (6,2) δεν είναι κενό, αλλά περιέχει '┴'.} \\
\\ 
\texttt{Έχετε συσσωρεύσει συλλογικά 18 από 16 ποινές.} \\
\texttt{Παρακαλώ δοκιμάστε ξανά!} \\
            }
        }
    }
    & \\ \\

    \theutterance \stepcounter{utterance}  
    & & & \multicolumn{4}{p{0.6\linewidth}}{
        \cellcolor[rgb]{0.9,0.9,0.9}{
            \makecell[{{p{\linewidth}}}]{
                \texttt{\tiny{[GM$\langle$P2]}}
                \texttt{ΚΙΝΗΣΗ: P, (7, 3)} \\
            }
        }
    }
     \\ \\

    \theutterance \stepcounter{utterance}  
    & & & \multicolumn{2}{p{0.3\linewidth}}{
        \cellcolor[rgb]{0.9,0.9,0.9}{
            \makecell[{{p{\linewidth}}}]{
                \texttt{\tiny{[GM$|$GM]}}
                \texttt{ΚΙΝΗΣΗ: P, (7, 3)} \\
            }
        }
    }
    & & \\ \\

    \theutterance \stepcounter{utterance}  
    & & & \multicolumn{2}{p{0.3\linewidth}}{
        \cellcolor[rgb]{0.9,0.9,0.9}{
            \makecell[{{p{\linewidth}}}]{
                \texttt{\tiny{[GM$|$GM]}}
                \texttt{Ποινή: Το (7,3) δεν είναι κενό, αλλά περιέχει 'P'.} \\
\\ 
\texttt{Έχετε συσσωρεύσει συλλογικά 19 από 16 ποινές.} \\
\texttt{Παρακαλώ δοκιμάστε ξανά!} \\
            }
        }
    }
    & & \\ \\

    \theutterance \stepcounter{utterance}  
    & & \multicolumn{4}{p{0.6\linewidth}}{
        \cellcolor[rgb]{0.9,0.9,0.9}{
            \makecell[{{p{\linewidth}}}]{
                \texttt{\tiny{[GM$\rangle$P2]}}
                \texttt{Ποινή: Το (7,3) δεν είναι κενό, αλλά περιέχει 'P'.} \\
\\ 
\texttt{Έχετε συσσωρεύσει συλλογικά 19 από 16 ποινές.} \\
\texttt{Παρακαλώ δοκιμάστε ξανά!} \\
            }
        }
    }
    & \\ \\

    \theutterance \stepcounter{utterance}  
    & & & \multicolumn{4}{p{0.6\linewidth}}{
        \cellcolor[rgb]{0.9,0.9,0.9}{
            \makecell[{{p{\linewidth}}}]{
                \texttt{\tiny{[GM$\langle$P2]}}
                \texttt{ΚΙΝΗΣΗ: P, (7, 2)} \\
            }
        }
    }
     \\ \\

    \theutterance \stepcounter{utterance}  
    & & & \multicolumn{2}{p{0.3\linewidth}}{
        \cellcolor[rgb]{0.9,0.9,0.9}{
            \makecell[{{p{\linewidth}}}]{
                \texttt{\tiny{[GM$|$GM]}}
                \texttt{ΚΙΝΗΣΗ: P, (7, 2)} \\
            }
        }
    }
    & & \\ \\

    \theutterance \stepcounter{utterance}  
    & & & \multicolumn{2}{p{0.3\linewidth}}{
        \cellcolor[rgb]{0.9,0.9,0.9}{
            \makecell[{{p{\linewidth}}}]{
                \texttt{\tiny{[GM$|$GM]}}
                \texttt{Ποινή: Το (7,2) δεν είναι κενό, αλλά περιέχει '─'.} \\
\\ 
\texttt{Έχετε συσσωρεύσει συλλογικά 20 από 16 ποινές.} \\
\texttt{Παρακαλώ δοκιμάστε ξανά!} \\
            }
        }
    }
    & & \\ \\

    \theutterance \stepcounter{utterance}  
    & & \multicolumn{4}{p{0.6\linewidth}}{
        \cellcolor[rgb]{0.9,0.9,0.9}{
            \makecell[{{p{\linewidth}}}]{
                \texttt{\tiny{[GM$\rangle$P2]}}
                \texttt{Ποινή: Το (7,2) δεν είναι κενό, αλλά περιέχει '─'.} \\
\\ 
\texttt{Έχετε συσσωρεύσει συλλογικά 20 από 16 ποινές.} \\
\texttt{Παρακαλώ δοκιμάστε ξανά!} \\
            }
        }
    }
    & \\ \\

    \theutterance \stepcounter{utterance}  
    & & & \multicolumn{4}{p{0.6\linewidth}}{
        \cellcolor[rgb]{0.9,0.9,0.9}{
            \makecell[{{p{\linewidth}}}]{
                \texttt{\tiny{[GM$\langle$P2]}}
                \texttt{ΚΙΝΗΣΗ: P, (5, 2)} \\
            }
        }
    }
     \\ \\

    \theutterance \stepcounter{utterance}  
    & & & \multicolumn{2}{p{0.3\linewidth}}{
        \cellcolor[rgb]{0.9,0.9,0.9}{
            \makecell[{{p{\linewidth}}}]{
                \texttt{\tiny{[GM$|$GM]}}
                \texttt{ΚΙΝΗΣΗ: P, (5, 2)} \\
            }
        }
    }
    & & \\ \\

    \theutterance \stepcounter{utterance}  
    & & & \multicolumn{2}{p{0.3\linewidth}}{
        \cellcolor[rgb]{0.9,0.9,0.9}{
            \makecell[{{p{\linewidth}}}]{
                \texttt{\tiny{[GM$|$GM]}}
                \texttt{Ποινή: Το (5,2) δεν είναι κενό, αλλά περιέχει '─'.} \\
\\ 
\texttt{Έχετε συσσωρεύσει συλλογικά 21 από 16 ποινές.} \\
\texttt{Παρακαλώ δοκιμάστε ξανά!} \\
            }
        }
    }
    & & \\ \\

    \theutterance \stepcounter{utterance}  
    & & \multicolumn{4}{p{0.6\linewidth}}{
        \cellcolor[rgb]{0.9,0.9,0.9}{
            \makecell[{{p{\linewidth}}}]{
                \texttt{\tiny{[GM$\rangle$P2]}}
                \texttt{Ποινή: Το (5,2) δεν είναι κενό, αλλά περιέχει '─'.} \\
\\ 
\texttt{Έχετε συσσωρεύσει συλλογικά 21 από 16 ποινές.} \\
\texttt{Παρακαλώ δοκιμάστε ξανά!} \\
            }
        }
    }
    & \\ \\

    \theutterance \stepcounter{utterance}  
    & & & \multicolumn{4}{p{0.6\linewidth}}{
        \cellcolor[rgb]{0.9,0.9,0.9}{
            \makecell[{{p{\linewidth}}}]{
                \texttt{\tiny{[GM$\langle$P2]}}
                \texttt{ΚΙΝΗΣΗ: P, (6, 3)} \\
            }
        }
    }
     \\ \\

    \theutterance \stepcounter{utterance}  
    & & & \multicolumn{2}{p{0.3\linewidth}}{
        \cellcolor[rgb]{0.9,0.9,0.9}{
            \makecell[{{p{\linewidth}}}]{
                \texttt{\tiny{[GM$|$GM]}}
                \texttt{ΚΙΝΗΣΗ: P, (6, 3)} \\
            }
        }
    }
    & & \\ \\

    \theutterance \stepcounter{utterance}  
    & & & \multicolumn{2}{p{0.3\linewidth}}{
        \cellcolor[rgb]{0.9,0.9,0.9}{
            \makecell[{{p{\linewidth}}}]{
                \texttt{\tiny{[GM$|$GM]}}
                \texttt{P moved to (6, 3)} \\
            }
        }
    }
    & & \\ \\

    \theutterance \stepcounter{utterance}  
    & & \multicolumn{4}{p{0.6\linewidth}}{
        \cellcolor[rgb]{0.9,0.9,0.9}{
            \makecell[{{p{\linewidth}}}]{
                \texttt{\tiny{[P1$\langle$GM]}}
                \texttt{Μεταφέρθηκε επιτυχώς το 'A' στο (5,3). Το ενημερωμένο πλέγμα σας μοιάζει ως εξής:} \\
\texttt{```} \\
\texttt{ 1234567} \\
\texttt{╔══╤══╤═╗} \\
\texttt{║PM│SO│◌║ 1} \\
\texttt{╟──┼──┴─╢ 2} \\
\texttt{║◌◌│◌A◌◌║ 3} \\
\texttt{╟──┘◌◌◌◌║ 4} \\
\texttt{║D◌◌◌◌◌◌║ 5} \\
\texttt{║◌◌┌─┬──╢ 6} \\
\texttt{║T◌│◌│◌◌║ 7} \\
\texttt{╚══╧═╧══╝} \\
\texttt{```} \\
\texttt{Αυτή τη στιγμή παίζετε τον γύρο 7 από συνολικά 28.} \\
\texttt{Έχετε συσσωρεύσει συλλογικά 21 από 16 ποινές.} \\
\texttt{Ο άλλος παίκτης μετακίνησε ένα αντικείμενο στο πλέγμα του.} \\
\\ 
\\ 
\texttt{Ποια είναι η επόμενη εντολή σας;} \\
            }
        }
    }
    & \\ \\

    \theutterance \stepcounter{utterance}  
    & \multicolumn{4}{p{0.6\linewidth}}{
        \cellcolor[rgb]{0.9,0.9,0.9}{
            \makecell[{{p{\linewidth}}}]{
                \texttt{\tiny{[P1$\rangle$GM]}}
                \texttt{ΠΕΣ: Συγγνώμη για τα λάθη, έχουμε ξεπεράσει το όριο ποινών και το παιχνίδι τελείωσε.} \\
            }
        }
    }
    & & \\ \\

    \theutterance \stepcounter{utterance}  
    & & & \multicolumn{2}{p{0.3\linewidth}}{
        \cellcolor[rgb]{0.9,0.9,0.9}{
            \makecell[{{p{\linewidth}}}]{
                \texttt{\tiny{[GM$|$GM]}}
                \texttt{ΠΕΣ: Συγγνώμη για τα λάθη, έχουμε ξεπεράσει το όριο ποινών και το παιχνίδι τελείωσε.} \\
            }
        }
    }
    & & \\ \\

    \theutterance \stepcounter{utterance}  
    & & \multicolumn{4}{p{0.6\linewidth}}{
        \cellcolor[rgb]{0.9,0.9,0.9}{
            \makecell[{{p{\linewidth}}}]{
                \texttt{\tiny{[GM$\rangle$P2]}}
                \texttt{Μεταφέρθηκε επιτυχώς το 'P' στο (6,3). Το ενημερωμένο πλέγμα σας μοιάζει ως εξής:} \\
\texttt{```} \\
\texttt{ 1234567} \\
\texttt{╔══╤══╤═╗} \\
\texttt{║DM│SA│T║ 1} \\
\texttt{╟──┼──┴─╢ 2} \\
\texttt{║◌◌│◌◌P◌║ 3} \\
\texttt{╟──┘◌◌◌O║ 4} \\
\texttt{║◌◌◌◌◌◌◌║ 5} \\
\texttt{║◌◌┌─┬──╢ 6} \\
\texttt{║◌◌│◌│◌◌║ 7} \\
\texttt{╚══╧═╧══╝} \\
\texttt{```} \\
\texttt{Αυτή τη στιγμή παίζετε τον γύρο 8 από συνολικά 28.} \\
\texttt{Έχετε συσσωρεύσει συλλογικά 21 από 16 ποινές.} \\
\texttt{Ο άλλος παίκτης έστειλε αυτό το μήνυμα:} \\
\\ 
\\ 
\texttt{"Συγγνώμη για τα λάθη, έχουμε ξεπεράσει το όριο ποινών και το παιχνίδι τελείωσε."} \\
\\ 
\\ 
\texttt{Ποια είναι η επόμενη εντολή σας;} \\
            }
        }
    }
    & \\ \\

    \theutterance \stepcounter{utterance}  
    & & & \multicolumn{4}{p{0.6\linewidth}}{
        \cellcolor[rgb]{0.9,0.9,0.9}{
            \makecell[{{p{\linewidth}}}]{
                \texttt{\tiny{[GM$\langle$P2]}}
                \texttt{ΠΕΣ: τελείωσα!} \\
            }
        }
    }
     \\ \\

    \theutterance \stepcounter{utterance}  
    & & & \multicolumn{2}{p{0.3\linewidth}}{
        \cellcolor[rgb]{0.9,0.9,0.9}{
            \makecell[{{p{\linewidth}}}]{
                \texttt{\tiny{[GM$|$GM]}}
                \texttt{ΠΕΣ: τελείωσα!} \\
            }
        }
    }
    & & \\ \\

    \theutterance \stepcounter{utterance}  
    & & \multicolumn{4}{p{0.6\linewidth}}{
        \cellcolor[rgb]{0.9,0.9,0.9}{
            \makecell[{{p{\linewidth}}}]{
                \texttt{\tiny{[P1$\langle$GM]}}
                \texttt{Το μήνυμά σας μεταφέρθηκε στον άλλο παίκτη.} \\
\\ 
\\ 
\texttt{Αυτή τη στιγμή παίζετε τον γύρο 8 από συνολικά 28.} \\
\texttt{Έχετε συσσωρεύσει συλλογικά 21 από 16 ποινές.} \\
\texttt{Ο άλλος παίκτης έστειλε αυτό το μήνυμα:} \\
\\ 
\\ 
\texttt{"τελείωσα!"} \\
\\ 
\\ 
\texttt{Ποια είναι η επόμενη εντολή σας;} \\
            }
        }
    }
    & \\ \\

    \theutterance \stepcounter{utterance}  
    & \multicolumn{4}{p{0.6\linewidth}}{
        \cellcolor[rgb]{0.9,0.9,0.9}{
            \makecell[{{p{\linewidth}}}]{
                \texttt{\tiny{[P1$\rangle$GM]}}
                \texttt{ΠΕΣ: τελείωσα!} \\
            }
        }
    }
    & & \\ \\

    \theutterance \stepcounter{utterance}  
    & & & \multicolumn{2}{p{0.3\linewidth}}{
        \cellcolor[rgb]{0.9,0.9,0.9}{
            \makecell[{{p{\linewidth}}}]{
                \texttt{\tiny{[GM$|$GM]}}
                \texttt{ΠΕΣ: τελείωσα!} \\
            }
        }
    }
    & & \\ \\

    \theutterance \stepcounter{utterance}  
    & & & \multicolumn{2}{p{0.3\linewidth}}{
        \cellcolor[rgb]{0.9,0.9,0.9}{
            \makecell[{{p{\linewidth}}}]{
                \texttt{\tiny{[GM$|$GM]}}
                \texttt{true} \\
            }
        }
    }
    & & \\ \\

    \theutterance \stepcounter{utterance}  
    & & & \multicolumn{2}{p{0.3\linewidth}}{
        \cellcolor[rgb]{0.9,0.9,0.9}{
            \makecell[{{p{\linewidth}}}]{
                \texttt{\tiny{[GM$|$GM]}}
                \texttt{Initial states:} \\
\texttt{```} \\
\texttt{Player 1:} \\
\texttt{ 1234567} \\
\texttt{╔══╤══╤═╗} \\
\texttt{║◌P│◌M│◌║ 1} \\
\texttt{╟──┼──┴─╢ 2} \\
\texttt{║◌◌│◌S◌O║ 3} \\
\texttt{╟──┘◌A◌◌║ 4} \\
\texttt{║D◌◌◌◌◌◌║ 5} \\
\texttt{║◌◌┌─┬──╢ 6} \\
\texttt{║T◌│◌│◌◌║ 7} \\
\texttt{╚══╧═╧══╝} \\
\\ 
\\ 
\texttt{Player 2:} \\
\texttt{ 1234567} \\
\texttt{╔══╤══╤═╗} \\
\texttt{║◌◌│◌◌│T║ 1} \\
\texttt{╟──┼──┴─╢ 2} \\
\texttt{║◌◌│◌S◌P║ 3} \\
\texttt{╟──┘◌◌MO║ 4} \\
\texttt{║D◌◌◌◌◌◌║ 5} \\
\texttt{║◌◌┌─┬──╢ 6} \\
\texttt{║◌◌│A│◌◌║ 7} \\
\texttt{╚══╧═╧══╝} \\
\\ 
\texttt{```} \\
            }
        }
    }
    & & \\ \\

    \theutterance \stepcounter{utterance}  
    & & & \multicolumn{2}{p{0.3\linewidth}}{
        \cellcolor[rgb]{0.9,0.9,0.9}{
            \makecell[{{p{\linewidth}}}]{
                \texttt{\tiny{[GM$|$GM]}}
                \texttt{End states:} \\
\texttt{```} \\
\texttt{Player 1:} \\
\texttt{ 1234567} \\
\texttt{╔══╤══╤═╗} \\
\texttt{║PM│SO│◌║ 1} \\
\texttt{╟──┼──┴─╢ 2} \\
\texttt{║◌◌│◌A◌◌║ 3} \\
\texttt{╟──┘◌◌◌◌║ 4} \\
\texttt{║D◌◌◌◌◌◌║ 5} \\
\texttt{║◌◌┌─┬──╢ 6} \\
\texttt{║T◌│◌│◌◌║ 7} \\
\texttt{╚══╧═╧══╝} \\
\\ 
\\ 
\texttt{Player 2:} \\
\texttt{ 1234567} \\
\texttt{╔══╤══╤═╗} \\
\texttt{║DM│SA│T║ 1} \\
\texttt{╟──┼──┴─╢ 2} \\
\texttt{║◌◌│◌◌P◌║ 3} \\
\texttt{╟──┘◌◌◌O║ 4} \\
\texttt{║◌◌◌◌◌◌◌║ 5} \\
\texttt{║◌◌┌─┬──╢ 6} \\
\texttt{║◌◌│◌│◌◌║ 7} \\
\texttt{╚══╧═╧══╝} \\
\\ 
\texttt{```} \\
            }
        }
    }
    & & \\ \\

    \theutterance \stepcounter{utterance}  
    & & & \multicolumn{2}{p{0.3\linewidth}}{
        \cellcolor[rgb]{0.9,0.9,0.9}{
            \makecell[{{p{\linewidth}}}]{
                \texttt{\tiny{[GM$|$GM]}}
                \texttt{* success: True} \\
\texttt{* lose: False} \\
\texttt{* aborted: False} \\
\texttt{{-}{-}{-}{-}{-}{-}{-}} \\
\texttt{* Shifts: 9.00} \\
\texttt{* Max Shifts: 12.00} \\
\texttt{* Min Shifts: 6.00} \\
\texttt{* End Distance Sum: 23.48} \\
\texttt{* Init Distance Sum: 21.20} \\
\texttt{* Expected Distance Sum: 29.33} \\
\texttt{* Penalties: 21.00} \\
\texttt{* Max Penalties: 16.00} \\
\texttt{* Rounds: 8.00} \\
\texttt{* Max Rounds: 28.00} \\
\texttt{* Object Count: 7.00} \\
            }
        }
    }
    & & \\ \\

\end{supertabular}
}

\end{document}
